% SPDX-FileCopyrightText: 2026 Jonas Whidden
% SPDX-License-Identifier: Apache-2.0 OR CC-BY-4.0

\documentclass[11pt]{amsart}

\usepackage[T1]{fontenc}
\usepackage{lmodern}
\usepackage[nopatch=footnote]{microtype}
\usepackage[margin=1.1in]{geometry}
\usepackage{amsmath,amssymb,amsthm,mathtools}
\usepackage{array,booktabs,longtable}
\usepackage{enumitem}
\usepackage{xcolor}
\usepackage{hyperref}

\definecolor{linkblue}{HTML}{174A7E}
\hypersetup{
  colorlinks=true,
  linkcolor=linkblue,
  citecolor=linkblue,
  urlcolor=linkblue,
  pdftitle={Robin's 1984 criterion for the Riemann hypothesis},
  pdfauthor={Jonas Whidden}
}

\newtheorem{theorem}{Theorem}[section]
\newtheorem{proposition}[theorem]{Proposition}
\newtheorem{lemma}[theorem]{Lemma}
\newtheorem{corollary}[theorem]{Corollary}
\theoremstyle{definition}
\newtheorem{definition}[theorem]{Definition}
\theoremstyle{remark}
\newtheorem{remark}[theorem]{Remark}

\numberwithin{equation}{section}

\newcommand{\RH}{\mathrm{RH}}
\newcommand{\CA}{\mathrm{CA}}
\newcommand{\N}{\mathbb{N}}
\newcommand{\R}{\mathbb{R}}
\newcommand{\C}{\mathbb{C}}
\newcommand{\euler}{\gamma}
\newcommand{\sig}{\sigma}
\newcommand{\A}{\mathcal{A}}
\newcommand{\M}{\mathcal{M}}
\newcommand{\E}{\mathcal{E}}
\newcommand{\Pcost}{\mathcal{C}}
\newcommand{\Pprod}{\mathfrak{P}}
\newcommand{\eps}{\varepsilon}
\newcommand{\leanname}[1]{\texttt{#1}}
\newcommand{\repositoryurl}{https://github.com/kimihiro64/Robin1984}
\newcommand{\codefile}[1]{%
  \href{\repositoryurl/blob/main/#1}{\nolinkurl{#1}}}
\newcommand{\codedir}[1]{%
  \href{\repositoryurl/tree/main/#1}{\nolinkurl{#1}}}
\newcommand{\codelink}[2]{%
  \href{\repositoryurl/blob/main/#1}{\texttt{#2}}}
\newcommand{\dirlink}[2]{%
  \href{\repositoryurl/tree/main/#1}{\texttt{#2}}}
\newcommand{\commitlink}[2]{%
  \href{\repositoryurl/commit/#1}{\texttt{#2}}}
\newcommand{\runlink}[2]{%
  \href{\repositoryurl/actions/runs/#1}{#2}}
\newcommand{\pntrepositoryurl}{%
  https://github.com/kimihiro64/PrimeNumberTheoremAnd}
\newcommand{\pntrevision}{f8f58c749d6cde8a641348fcd5e4702993651cd6}
\newcommand{\pntfile}[2]{%
  \href{\pntrepositoryurl/blob/\pntrevision/#1}{\texttt{#2}}}
\newcommand{\actionsurl}{%
  https://github.com/kimihiro64/Robin1984/actions/workflows/ci.yml}
\newcommand{\releasesurl}{%
  https://github.com/kimihiro64/Robin1984/releases}

\title[Robin's 1984 criterion]{Robin's 1984 criterion for the Riemann hypothesis:\\
a formally verified proof}
\author{Jonas Whidden}
\date{September 2026}
\subjclass[2020]{11M26, 11N37, 03B35}
\keywords{Riemann hypothesis, sum of divisors, Robin's inequality,
colossally abundant numbers, Nicolas criterion, formal verification, Lean}

\begin{document}
\raggedbottom
\begin{abstract}
Guy Robin proved in 1984 that the Riemann hypothesis is equivalent to the
strict inequality
\[
  \sig(n)<e^{\euler}n\log\log n \qquad(n>5040).
\]
We give a detailed mathematical reconstruction paired with a formally verified
Lean~4 proof.  The argument reduces arbitrary integers to colossally abundant
extrema, proves the bounded range with exact rational certificates, establishes
the large-height implication under the Riemann hypothesis through a weighted
explicit formula and a strict coefficient gap, and proves the converse by
turning Nicolas--Landau oscillations into least-common-multiple counterexamples.
The source-specific weighted formula and oscillation argument are proved in
repository-owned modules, while substantial foundational analysis is imported
from pinned public libraries: axiom-free here does not mean dependency-free or
redeveloped from first principles.

A declaration-level map connects the exposition to the checked source.
Comparator verifies that the solution proves the exact Mathlib-only challenge
statements and respects the configured axiom boundary; Lean's kernel and the
independent NanoDa kernel then replay the export.  Public source, certificates,
CI logs, release artifacts, and reproduction instructions make this assurance
chain inspectable.  The reconstruction used a human-directed automated context
system to track sources, theorem state, dependencies, numerical evidence, and
verification status.  The paper documents its role through retained outputs
and two concrete failure-and-revision cases, without claiming a measured
productivity advantage.  Reusable outputs include a generic positive-transform
theorem, multiplicity-aware zero infrastructure, and a theorem-parametric
finite-certificate checker.  Robin's result remains an equivalence criterion;
it does not prove the Riemann hypothesis.
\end{abstract}
\maketitle

\begin{center}
\small
Copyright \textcopyright\ 2026 Jonas Whidden. This paper is available under
the \href{https://www.apache.org/licenses/LICENSE-2.0}{Apache License 2.0} or
\href{https://creativecommons.org/licenses/by/4.0/}{CC BY 4.0}, at your option.
\end{center}

\section{Introduction}

For a positive integer $n$, let
\[
  \sig(n)=\sum_{d\mid n}d,
  \qquad
  \A(n)=\frac{\sig(n)}{n},
\]
and let $\euler$ be the Euler--Mascheroni constant.  Robin's theorem is the
following exact criterion.  Throughout, $\RH$ denotes the standard
zero-locus formulation: every zero $s$ of the analytically continued
Riemann zeta function which is neither a negative even integer nor the pole
at $1$ satisfies $\Re s=1/2$.  This is also Mathlib's definition of
\leanname{RiemannHypothesis}; it is stronger in presentation, but equivalent
in the usual way, to saying that zeta has no zero with $\Re s>1/2$.

\begin{theorem}[Robin~\cite{Robin1984}]\label{thm:main}
The Riemann hypothesis holds if and only if
\begin{equation}\label{eq:robin}
  \sig(n)<e^{\euler}n\log\log n
\end{equation}
for every integer $n>5040$.
\end{theorem}

The cutoff is part of the theorem.  In particular, the result is not merely
an eventual estimate.  It combines an analytic argument above an explicit
height with a finite proof below that height.

Robin's original article~\cite{Robin1984} is in French.  Although a scanned
PDF is now available online, the paper remains comparatively difficult to
discover through ordinary bibliographic channels.  One purpose of the present
paper is therefore expository: to give an accessible English, proof-oriented
account of the argument while preserving the attribution and source trail of
the original.  Nicolas's related Euler-function criterion has a later English
treatment~\cite{Nicolas2012}; the present paper concerns Robin's distinct
sum-of-divisors argument.

The complete Lean source, certificate data, reproducibility instructions,
and release artifacts are available at
\begin{center}
  \url{\repositoryurl}.
\end{center}
Section~\ref{sec:formalization} links the public theorem boundary and the
principal modules for each mathematical stage, while the finite-range
section links the exact certificate rows directly.

The second principal result makes Robin's extremal reduction explicit.  For
$\eps>0$, put
\begin{equation}\label{eq:ca-objective}
  Q_{\eps}(m)=\frac{\sig(m)}{m^{1+\eps}}
             =\frac{\A(m)}{m^{\eps}}.
\end{equation}
We say that $N>1$ is colossally abundant with parameter $\eps$ if
$Q_{\eps}(m)\le Q_{\eps}(N)$ for every $m>1$.

\begin{theorem}\label{thm:ca-main}
The Riemann hypothesis holds if and only if \eqref{eq:robin} holds for every
$N>5040$ which is colossally abundant with some positive parameter.
\end{theorem}

Colossally abundant numbers were introduced by Alaoglu and
Erd\H{o}s~\cite{AlaogluErdos1944}.  The criterion itself and its use of those
extrema are due to Robin.  The converse uses the oscillation theorem of
Nicolas~\cite{Nicolas1983} together with Landau's positive-transform
principle~\cite{Landau1905}.  The explicit estimates for prime functions used
below trace to Rosser and Schoenfeld~\cite{RosserSchoenfeld1962} and to Costa
Pereira~\cite{CostaPereira1985,CostaPereira1987}.  Our purpose is not to claim
new priority for those ingredients, but to give a complete proof whose exact
logical and numerical dependencies have been checked.

The new contribution is consequently one of reconstruction, reusable formal
infrastructure, and verification.  The analytic argument is carried through
a multiplicity-counted zero divisor; Landau's positive-transform principle is
isolated as a generic theorem about moment-generating functions; the finite
range is reduced to a reusable soundness theorem plus exact data; and the
public biconditional is assembled without a project axiom.  This separation
is useful beyond the present theorem: source mathematics, generic analytic
lemmas, untrusted certificate discovery, kernel-checked certificate
validation, and final theorem assembly remain distinct layers.

The development was assisted by an automated research and formalization
context system.  It maintained source provenance, exact theorem targets,
definition and import dependencies, live proof obligations, failed routes,
numerical evidence, and verification status while human direction selected
the mathematical claims and accepted the final statements.  Section
\ref{sec:formalization} describes this architecture and its trust boundary.
The claims made about that system concern the retained source, certificates,
documented revisions, and validation records.  No controlled comparison of
human effort, elapsed development time, or error rate is claimed.

The proof has four parts.  Section~\ref{sec:extremal} reduces arbitrary
integers to colossally abundant numbers, the global maximizers of
$Q_\eps$.  Section~\ref{sec:finite} closes the bounded range.
Section~\ref{sec:rh-forward} proves the inequality under $\RH$ from a strict
explicit coefficient gap.  Sections~\ref{sec:oscillation} and~\ref{sec:lcm}
give the Nicolas--Landau converse: an off-critical zero forces recurring
prime-product oscillations, which least common multiples turn into integer
counterexamples.  At the highest level the two directions are
\[
 \RH\ \Longrightarrow\ \text{explicit prime-function bounds}\
 \Longrightarrow\ \M(n)<0,
\]
and
\[
 \neg\RH\ \Longrightarrow\ \text{negative Nicolas oscillations}\
 \Longrightarrow\ \M(L_P)>0\ \text{for arbitrarily large }P.
\]
The finite theorem joins the first chain to the exact cutoff $5040$, and the
LCM identity joins the second chain to literal integer counterexamples.  The
dedicated main-text account of the Lean organization and checking is confined
to Section~\ref{sec:formalization}; the appendices contain supporting audits.

\section{Roadmap and informal guide}\label{sec:informal}

This section gives the argument without its explicit constants or
formalization details.  It is intended as a map for readers who do not work
regularly with analytic number theory.  The later sections prove each step in
the form actually needed for the theorem.

The proof and its checked endpoints can be read as the following directed
graph.  The declaration names are printed here so that a reader can move from
the mathematical route to the exact Lean statement without searching the
repository.
\begin{center}
\small
\begin{tabular}{@{}>{\raggedright\arraybackslash}p{.25\textwidth}
  >{\raggedright\arraybackslash}p{.66\textwidth}@{}}
\toprule
mathematical stage & paper endpoint and checked declaration \\
\midrule
CA extremal reduction
& Proposition~\ref{prop:ca-reduction};
  \path{nativeRobinInequalityAll_iff_colossallyAbundant} in
  \codelink{Robin1984/ColossallyAbundant/CAReduction.lean}{CAReduction}
  \\
finite range
& Theorem~\ref{thm:finite}; \path{nativeRobinInequality_finite_log} in
  \codelink{Robin1984/Finite/FiniteComplete.lean}{FiniteComplete}
  \\
large-height RH branch
& Theorem~\ref{thm:rh-forward};
  \path{nativeRobinInequality_of_RH_large_log} in
  \codelink{Robin1984/Equivalence/LargeHeightRobin.lean}{LargeHeightRobin}
  \\
false-RH converse
& Theorem~\ref{thm:converse};
  \path{riemannHypothesis_of_nativeRobinInequalityAll} in
  \codelink{Robin1984/NicolasLandau/NicolasLandauRobinBridge.lean}{LCM bridge}
  \\
public biconditionals
& Theorems~\ref{thm:main} and~\ref{thm:ca-main};
  \path{Robin1984.robin_inequality_iff_riemannHypothesis} and
  \path{Robin1984.riemannHypothesis_iff_colossallyAbundant_robin} in
  \codelink{Solution.lean}{Solution}
  \\
\bottomrule
\end{tabular}
\end{center}

\subsection{What Robin's inequality measures}

The ratio $\A(n)=\sig(n)/n$ measures how rich $n$ is in divisors relative to
its size.  If
\[
 n=\prod_p p^{a_p},
\]
then $\A(n)$ is a product of local factors, one for each prime dividing $n$.
Small primes increase this ratio most efficiently, and raising a prime to a
higher power gives progressively smaller additional gains.  Integers built
from many small prime powers are therefore the difficult cases for an upper
bound on $\A(n)$.

Robin's inequality says that even these difficult integers satisfy
\[
 \A(n)<e^\euler\log\log n
\]
once $n>5040$, exactly when $\RH$ is true.  Taking logarithms turns the
comparison into the sign of
\[
 \M(n)=\log\A(n)-\euler-\log\log\log n.
\]
The proof is consequently a search for structural reasons why this margin is
negative, or, if $\RH$ fails, why it must sometimes be positive.

\subsection{Why the extremal integers suffice}

For a positive penalty $\eps$, the quantity
\[
 Q_\eps(m)=\A(m)m^{-\eps}
\]
rewards divisors but charges for size.  Its maximizers are the colossally
abundant integers.  They are useful because the penalty separates into
independent prime-power decisions: include a prime-power event when its
marginal gain exceeds its size cost, exclude it when the cost is larger.

Given a general integer $n$, choose the penalty
\[
 \eps=\frac{1}{\log n\,\log\log n}.
\]
Let $q$ maximize $Q_\eps$.  Maximality says that the logarithmic abundancy of
$n$ lies below the affine line determined by $q$ and this penalty.  The
function appearing in Robin's bound is concave, and the chosen penalty is its
tangent slope at $n$.  The tangent inequality therefore transfers Robin's
inequality from $q$ back to $n$.  A short exact calculation at the transition
$5040\leftrightarrow55440$ ensures that the chosen maximizer lies above the
exceptional cutoff whenever this reduction is used.  The bounded interval
not covered by that argument is settled separately.  This explains why
checking colossally abundant integers is both sufficient and necessary.

\subsection{Why the Riemann hypothesis implies the inequality}

Write $H=\log n$.  The prime-power factorization of $n$ expresses
$\log\A(n)$ as a sum of marginal gains.  Completing this sum to a convenient
set of small primes produces two competing quantities.  One is an analytic
error term, controlled through weighted estimates for the Chebyshev
functions $\vartheta$ and $\psi$.  The other is a definite arithmetic
deduction: an integer of height $H$ cannot take every small-prime gain for
free, so omitted primes or finite exponents impose a cost.

Under $\RH$, the weighted explicit formula bounds the analytic error on the
natural square-root scale
\[
 S(H)=H^{-1/2}(\log H)^{-1}.
\]
For $H\ge74500$, the proof shows that the total error is strictly less than
$(113/50)S(H)$, while a complete block of prime squares guarantees a cost of
at least $(113/50)S(H)$.  Hence error minus cost is negative, which is exactly
$\M(n)<0$.

The square-root scale has a simple spectral interpretation.  The explicit
formula rewrites the weighted prime-counting discrepancy as a superposition
of contributions resembling $H^\rho$, one for each nontrivial zero $\rho$ of
zeta.  Under $\RH$, every contribution has magnitude on the $H^{1/2}$ scale;
the weight damps the oscillation and the inverse-square zero sum controls the
aggregate amplitude.  The displayed error term is therefore a rigorous
envelope for all of those waves, not a numerical fit.

The prime cost is the complementary combinatorial picture.  An integer with
logarithmic size $H$ has a finite budget of prime exponents.  It cannot collect
every favorable small-prime contribution without spending that budget.  The
minimum in \eqref{eq:prime-cost} records whichever of two unavoidable losses
is cheaper for a prime.  The strict coefficient gap says that even this
conservative total loss is larger than the worst analytic envelope.

Below that analytic height, the argument is finite but not a decimal
computation.  Exact factorizations handle the first interval.  The remaining
range is covered by rational log-height rows.  Each row lists a finite box of
possible profitable prime powers and proves, by exact integer products and
rational logarithm bounds, that the Robin margin is negative throughout the
whole interval.  Thus the finite and analytic halves meet at $H=74500$
without an unverified numerical gap.

\subsection{Why the inequality implies the Riemann hypothesis}

The converse is best read contrapositively.  If $\RH$ is false, symmetry of
the zeta function supplies a nontrivial zero to the right of the critical
line.  That zero becomes a pole of a Mellin transform built from a weighted
prime-counting error.  Landau's positivity principle says that this transform
cannot remain analytic as far as that pole if the underlying error is
eventually of one sign.  Consequently the error has negative excursions of a
fixed power size arbitrarily far out.

Geometrically, the real axis of convergence acts like a wall.  Eventual
one-sidedness turns the Mellin transform, after $x=e^y$, into the Laplace
transform of a positive measure.  Positivity forces convergence to advance
whenever the transform analytically crosses that wall.  An off-critical zero,
however, creates a genuine pole just beyond it.  The two statements are
incompatible, so the assumed one-sided tail must fail repeatedly.  Section
\ref{sec:oscillation} replaces this picture by the exact continuation,
integrability, and pole calculations.

Nicolas's finite product $F(P)$ converts those analytic excursions into
negative excursions of $\log F(P)$.  To obtain actual integers, take
\[
 L_P=\operatorname{lcm}(1,2,\ldots,P).
\]
This least common multiple contains precisely the prime powers up to $P$ and
satisfies $\log L_P=\psi(P)$.  An exact identity expresses its Robin margin
as
\[
 \M(L_P)=-\log F(P)-D(P)-T(P),
\]
where the height loss $D(P)$ and the finite-prime-power reserve $T(P)$ are
nonnegative but smaller than the oscillation.  The negative excursions of
$\log F(P)$ therefore give positive values of $\M(L_P)$ for arbitrarily
large $P$.  These are violations of Robin's inequality above $5040$.
Hence universal validity of the inequality forces $\RH$.

The technical proof that follows adds the exact thresholds, constants,
convergence arguments, and finite data needed to make this outline rigorous;
it does not change the four-step logical route.

\section{Historical comparison and scope of novelty}\label{sec:comparison}

The theorem, the use of colossally abundant extrema, the weighted
prime-function estimates, and the Nicolas--Landau oscillation mechanism are
classical.  This paper claims no new unconditional theorem about the zeta
function and no priority for those ingredients.  Its contributions occur at
three different levels: an English reconstruction of Robin's proof; several
exact decompositions and interfaces chosen to make that proof formalizable;
and a checked Lean artifact with independent statement and kernel replay.
The distinction is summarized below.

\begin{center}
\small
\begin{tabular}{@{}>{\raggedright\arraybackslash}p{.18\textwidth}
  >{\raggedright\arraybackslash}p{.30\textwidth}
  >{\raggedright\arraybackslash}p{.42\textwidth}@{}}
\toprule
stage & classical source & present reconstruction \tabularnewline
\midrule
criterion and cutoff
& Robin's Theorem~1~\cite{Robin1984}
& The same theorem and cutoff, stated directly over Mathlib objects; no new
  mathematical claim. \tabularnewline
CA reduction
& Robin's Section~3, Proposition~1, compares successive CA integers using the
  convexity of $x\mapsto\eps x-\log\log x$.
& Proposition~\ref{prop:tangent-transfer} expresses the same supporting-line
  geometry at an arbitrary target and separates the exact
  $5040\leftrightarrow55440$ startup transition.  The reformulation and its
  quantifier-exact assembly are project work; the reduction is Robin's.
  \tabularnewline
RH implication
& Robin's Section~3, especially Lemmas~2--8 and Proposition~3, combines a
  weighted explicit formula with prime-product estimates, following Nicolas
  but requiring finer bounds.
& The complete exponent budget packages every actual exponent into an error
  term and a nonnegative prime cost, then proves a single exact coefficient
  gap at $H\ge74500$.  This is a project-authored formal decomposition and
  explicit cutoff, not a claim of a new asymptotic estimate or an optimal
  constant. \tabularnewline
finite range
& Robin reports computing 2347 CA integers below the analytic range and then
  checking the remaining integers below $55440$.
& Exact startup factorizations, a tangent interval, and rational certificate
  rows replace the historical computer calculation by kernel-checkable data
  plus a general soundness theorem. \tabularnewline
false-RH converse
& Robin invokes Nicolas's oscillation theorem, itself based on Landau's
  principle, and transfers the conclusion directly to CA integers.
& Sections~\ref{sec:oscillation} and~\ref{sec:lcm} reconstruct the transform
  argument, isolate the positive-transform theorem, and use least common
  multiples to obtain literal integer counterexamples before applying the CA
  equivalence.  This is an alternative formal proof architecture, not a new
  oscillation theorem attributed to this project. \tabularnewline
\bottomrule
\end{tabular}
\end{center}

Two formal interfaces merit a more precise explanation.  First, a classical
sum over zeta zeros is normally understood to repeat a zero according to its
analytic multiplicity.  That convention is not enough for the interchanges
needed in Lean.  The multiplicity-aware infrastructure supplies an explicit
index type for the divisor of $\xi$, a value map from indices to nonzero zeros,
inverse-norm summability, and the hypotheses licensing each exchange of an
integral with an infinite sum.  It proves no new fact about where the zeros
lie; it turns the classical convention into a reusable formal interface.

Second, the complete prime cost in \eqref{eq:prime-cost} is a
bookkeeping device.  Robin estimates several prime products separately.  The
present proof first sums the actual exponent choices of an arbitrary integer,
then completes to all relevant primes.  Each missing prime or prematurely
terminated exponent pays the smaller of two explicit penalties, so the
analytic error and arithmetic loss can be bounded in independent modules.
This modularity is useful for verification, but no theorem here says that it
is computationally minimal or analytically more efficient than every earlier
argument.

Third, Lemma~\ref{lem:nicolas-comparison} is called a ``finite Nicolas
comparison'' because it connects Nicolas's finite product to the analytic
error used in his oscillation theorem.  The product, Mertens series, and
partial-summation idea are classical.  The repository-owned contribution is
the exact identity
$\log F=R_M+R_\vartheta+K$, the explicit $4/x$ tail, and the separately proved
signs $R_\vartheta\le0$ and $K\le J$ that make the comparison usable without
asymptotic shorthand.  This is again a formal reconstruction and reusable
interface, not a new oscillation phenomenon.

\subsection{Relation to prior formal work}

The adjacent formal literature concerns the analytic objects around the
criterion rather than Robin's divisor-sum biconditional itself.
Gomes and Kontorovich~\cite{GomesKontorovich2020} formalized an equivalent
zero-locus formulation of the Riemann hypothesis through the Dirichlet eta
function and supplied a Mathlib implementation of its analytic prerequisites;
their development does not state Robin's inequality.  Loeffler and
Stoll~\cite{LoefflerStoll2025} describe the Mathlib development of zeta and
$L$-functions, including the formal statement of the Riemann hypothesis used
here.  The PrimeNumberTheoremAnd project~\cite{PrimeNumberTheoremAnd} supplies
Hadamard-factorization, completed-zeta, Mertens, and explicit prime-function
infrastructure on which the present reconstruction builds.

Those developments establish definitions and reusable analytic foundations.
The present artifact connects them to a different endpoint: Robin's exact
sum-of-divisors criterion, including the weighted endpoint, complete exponent
budget, finite certificate cover, Nicolas--Landau continuation, and LCM
transfer.  The author is not aware of an earlier mechanically checked proof of
this complete biconditional, but no exhaustive priority claim is made.  The
comparison asserted here is narrower and inspectable: none of the cited formal
developments contains the two public declarations proved by this repository.

The primary intended audience is therefore formal mathematics and analytic
number theory.  For AI-assisted mathematics, the contribution is a concrete
case study in maintaining fixed theorem targets, provenance, exact
certificates, and independent acceptance gates; it is not a new general
machine-learning model.  Section~\ref{sec:formalization} makes that case study
and its evidence explicit.

\section{Arithmetic preliminaries}

For $n>e^e$, set
\begin{equation}\label{eq:margin}
  \M(n)=\log \A(n)-\euler-\log\log\log n.
\end{equation}
All logarithms in \eqref{eq:margin} are then defined, and Robin's inequality
is equivalent to $\M(n)<0$.

Write $n=\prod_p p^{a_p}$.  For a prime $p$ and an integer $j\ge1$, define
the marginal gain and threshold
\begin{equation}\label{eq:event}
  g_{p,j}
  =\log\!\left(\frac{1-p^{-(j+1)}}{1-p^{-j}}\right),
  \qquad
  \tau_{p,j}=\frac{g_{p,j}}{\log p}.
\end{equation}
The event $(p,j)$ records the change from exponent $j-1$ to exponent $j$.

\begin{lemma}[Prime-power event decomposition]\label{lem:event-decomp}
For every positive integer $n$ and every $\eps>0$,
\begin{align}
  \log \A(n)
    &=\sum_p\sum_{1\le j\le a_p}g_{p,j},\label{eq:event-abundancy}\\
  \log Q_{\eps}(n)
    &=\sum_p\sum_{1\le j\le a_p}
       \bigl(g_{p,j}-\eps\log p\bigr).\label{eq:event-ca}
\end{align}
For fixed $p$, the sequences $g_{p,j}$ and $\tau_{p,j}$ are nonincreasing in
$j$.
\end{lemma}

\begin{proof}
The geometric-series identity gives
\[
  \frac{\sig(p^a)}{p^a}
  =\frac{1-p^{-(a+1)}}{1-p^{-1}}
  =\prod_{j=1}^{a}
    \frac{1-p^{-(j+1)}}{1-p^{-j}}.
\]
Multiplicativity and logarithms give \eqref{eq:event-abundancy}; subtracting
$\eps\log n$ gives \eqref{eq:event-ca}.  If $q=p^{-1}\in(0,1)$, direct
cross-multiplication shows that
$(1-q^{j+2})/(1-q^{j+1})\le(1-q^{j+1})/(1-q^j)$.
Taking logarithms and then dividing by $\log p>0$ proves the monotonicity.
\end{proof}

The reduced weight of an event at parameter $\eps$ is
\begin{equation}\label{eq:reduced-weight}
  w_{\eps}(p,j)=g_{p,j}-\eps\log p
               =(\tau_{p,j}-\eps)\log p.
\end{equation}
Thus events above the threshold have nonnegative reduced weight and events
below it have nonpositive reduced weight.  This is the elementary
complementary-slackness description of a colossally abundant factorization.

\begin{lemma}[Existence of a maximizer]\label{lem:ca-exists}
For every $\eps>0$ there exists an integer $N>1$ maximizing $Q_{\eps}$.
\end{lemma}

\begin{proof}
The divisor bound
\[
  \A(n)=\sum_{d\mid n}\frac1d
  \le \sum_{1\le d\le n}\frac1d
  \le 1+\log n
\]
implies
$Q_{\eps}(n)\le(1+\log n)n^{-\eps}\to0$.  Beyond a finite point the
objective is below its value at $2$.  A maximum over the remaining finite
set is therefore a global maximum.
\end{proof}

\section{The colossally abundant tangent reduction}\label{sec:extremal}

Put $H_n=\log n$ and
\begin{equation}\label{eq:lambda}
  \lambda(n)=\frac1{H_n\log H_n}.
\end{equation}
This is the derivative at $H_n$ of the concave function
$\phi(H)=\log\log H$.

\begin{lemma}[Global tangent inequality]\label{lem:tangent}
If $L,U>1$, then
\begin{equation}\label{eq:tangent}
  \log\log U-\log\log L
  \le \frac{U-L}{L\log L}.
\end{equation}
\end{lemma}

\begin{proof}
Apply $\log y\le y-1$ first to
$y=(\log U)/(\log L)$ and then to $y=U/L$.  Division by the positive
quantities $L$ and $\log L$ gives \eqref{eq:tangent}.  This proof applies
whether $U$ lies to the left or to the right of $L$.
\end{proof}

\begin{proposition}[Tangent transfer]\label{prop:tangent-transfer}
Let $n,q>5040$.  Suppose that
\begin{equation}\label{eq:tangent-objective}
  \log Q_{\lambda(n)}(n)
  \le \log Q_{\lambda(n)}(q).
\end{equation}
If Robin's inequality holds at $q$, then it holds at $n$.
\end{proposition}

\begin{proof}
Writing $L=H_n$ and $U=H_q$, inequality
\eqref{eq:tangent-objective} becomes
\[
  \log\A(n)-\lambda(n)L
  \le \log\A(q)-\lambda(n)U.
\]
By Lemma~\ref{lem:tangent},
$\phi(U)-\phi(L)\le\lambda(n)(U-L)$.  Adding the two inequalities yields
\[
  \log\A(n)-\phi(L)
  \le \log\A(q)-\phi(U).
\]
After subtracting $\euler$, the left- and right-hand sides are precisely
$\M(n)$ and $\M(q)$.  Hence $\M(q)<0$ implies $\M(n)<0$.
\end{proof}

The remaining point is to ensure that a maximizer at the tangent parameter
is itself above the exceptional cutoff.  The exact startup transition is
particularly useful.  Let
\begin{equation}\label{eq:epsilon-zero}
  \eps_0=\frac{\log(12/11)}{\log 11}.
\end{equation}
The event $(11,1)$ changes $5040$ into $55440=11\cdot5040$ and has threshold
$\eps_0$.

\begin{lemma}[The startup tie]\label{lem:startup-tie}
Both $5040$ and $55440$ maximize $Q_{\eps_0}$.  Moreover
\begin{equation}\label{eq:epsilon-enclosure}
  \frac1{34}<\eps_0<\frac1{22}.
\end{equation}
\end{lemma}

\begin{proof}
The factorization $5040=2^4 3^2 5\cdot7$ determines its included and first
excluded events.  Put
\[
  r_{p,j}=\frac{1-p^{-(j+1)}}{1-p^{-j}},
  \qquad \tau_{p,j}=\frac{\log r_{p,j}}{\log p}.
\]
The complete boundary calculation is small enough to display:
\begin{center}
\begin{tabular}{@{}ccl@{}}
\toprule
event & $r_{p,j}$ & comparison with $\eps_0$ \\
\midrule
$(2,4)$ & $31/30$ & $\tau_{2,4}>\eps_0$ \\
$(3,2)$ & $13/12$ & $\tau_{3,2}>\eps_0$ \\
$(5,1)$ & $6/5$   & $\tau_{5,1}>\eps_0$ \\
$(7,1)$ & $8/7$   & $\tau_{7,1}>\eps_0$ \\
\midrule
$(2,5)$ & $63/62$ & $\tau_{2,5}<\eps_0$ \\
$(3,3)$ & $40/39$ & $\tau_{3,3}<\eps_0$ \\
$(5,2)$ & $31/30$ & $\tau_{5,2}<\eps_0$ \\
$(7,2)$ & $57/56$ & $\tau_{7,2}<\eps_0$ \\
$(11,1)$ & $12/11$ & $\tau_{11,1}=\eps_0$ \\
\bottomrule
\end{tabular}
\end{center}
Here is the complete arithmetic behind the strict signs.  For an event with
base $p$ and gain ratio $r$, positivity of logarithms gives
\[
 p<r^N\ \Longrightarrow\ \frac1N<\frac{\log r}{\log p},
 \qquad
 r^N<p\ \Longrightarrow\ \frac{\log r}{\log p}<\frac1N.
\]
The cutoff enclosure~\eqref{eq:epsilon-enclosure} follows from the two exact
rational inequalities
\[
 11<(12/11)^{34},\qquad (12/11)^{22}<11.
\]
For the four included top events the kernel checks
\[
 2<(31/30)^{22},\quad 3<(13/12)^{22},\quad
 5<(6/5)^{22},\quad 7<(8/7)^{22};
\]
hence their thresholds exceed $1/22>\eps_0$.  For the four first-missing
events it checks
\[
 (63/62)^{34}<2,\quad (40/39)^{34}<3,\quad
 (31/30)^{34}<5,\quad (57/56)^{34}<7;
\]
hence their thresholds lie below $1/34<\eps_0$.  These are integer
comparisons after denominators are cleared; the tie $(11,1)$ is definitional.

For a fixed prime, monotonicity in $j$ propagates the displayed top-event
and first-missing-event signs through every layer.  Every prime not dividing
$5040$ is at least $11$, and the first-layer threshold is nonincreasing in
the prime; hence its reduced weight is nonpositive, with equality only at
$(11,1)$.  Equation~\eqref{eq:reduced-weight} now compares the event packet
of $5040$ term by term with every competitor.  Thus $5040$ maximizes
$Q_{\eps_0}$.  The tied event has weight zero, and adjoining it multiplies
$5040$ by $11$, proving the same maximality for $55440$.
\end{proof}

\begin{lemma}\label{lem:maximizer-above}
If $n\ge720720$ and $q$ maximizes $Q_{\lambda(n)}$, then $q>5040$.
\end{lemma}

\begin{proof}
The function $\lambda(n)$ is nonincreasing above $5040$.  Exact logarithmic
bounds give
\[
  \lambda(720720)<\frac1{34}<\eps_0,
\]
and hence $\lambda(n)<\eps_0$.  Write $\lambda=\lambda(n)$.  If
$q\le5040$, then maximality at $\eps_0$ gives
\[
 \log Q_\lambda(q)-\log Q_\lambda(5040)
 =\log Q_{\eps_0}(q)-\log Q_{\eps_0}(5040)
  +(\eps_0-\lambda)\log\frac q{5040}\le0.
\]
On the other hand, the exact tie of Lemma~\ref{lem:startup-tie} gives
\[
 \log Q_\lambda(55440)-\log Q_\lambda(5040)
   =(\eps_0-\lambda)\log 11>0.
\]
Thus no $q\le5040$ can maximize $Q_\lambda$, because $55440$ has strictly
larger objective value.
\end{proof}

\begin{proposition}[Reduction to colossally abundant integers]
\label{prop:ca-reduction}
Assume Robin's inequality for $5040<n<720720$.  Then the following are
equivalent:
\begin{enumerate}[label=(\roman*)]
  \item Robin's inequality holds for every $n>5040$;
  \item it holds for every colossally abundant $N>5040$.
\end{enumerate}
\end{proposition}

\begin{proof}
Only (ii)$\Rightarrow$(i) needs proof.  Let $n\ge720720$.  By
Lemma~\ref{lem:ca-exists}, choose a maximizer $q$ for
$Q_{\lambda(n)}$.  Lemma~\ref{lem:maximizer-above} gives $q>5040$, so (ii)
gives Robin's inequality at $q$.  Maximality gives
\eqref{eq:tangent-objective}, and Proposition~\ref{prop:tangent-transfer}
transfers the inequality to $n$.  The assumed finite interval handles the
remaining integers.
\end{proof}

\section{The finite range}\label{sec:finite}

The analytic argument in Section~\ref{sec:rh-forward} starts at
$\log n=74500$.  The exact cutoff therefore requires a proof for all
$5040<n$ below that logarithmic height.  We divide it into three pieces.

\begin{proposition}\label{prop:small-factor}
For $5041\le n<7560$,
\begin{equation}\label{eq:127-35}
  35\sig(n)\le127n.
\end{equation}
Consequently Robin's inequality holds throughout this interval.
\end{proposition}

\begin{proof}
For each integer in the interval, an exact prime-power factorization is
given.  Pairwise coprimality and the identity
$\sig(p^a)=1+p+\cdots+p^a$ reduce \eqref{eq:127-35} to natural-number
arithmetic.  To pass from \eqref{eq:127-35} to Robin's inequality, use the
exact bounds
\[
  e^{\euler}>\frac{177}{100},
  \qquad
  \log\log n>\frac{103}{50},
  \qquad
  \frac{127}{35}<\frac{177}{100}\frac{103}{50}.
\]
\end{proof}

\begin{proposition}\label{prop:startup-tangent}
Robin's inequality holds for $7560\le n<720720$.
\end{proposition}

\begin{proof}
The maximality of $5040$ at $\eps_0$ gives
\[
  \log\A(n)-\eps_0\log n
  \le \log\A(5040)-\eps_0\log5040.
\]
Set $y=\log n$.  It is enough to prove positivity of the explicit gap
\[
 G(y)=\euler+\log\log y-\log\A(5040)
      -\eps_0\bigl(y-\log5040\bigr),
\]
because the preceding inequality then places $\log\A(n)$ strictly below
$\euler+\log\log y$.  The curvature used at this step is explicit:
\[
 G''(y)=-\frac{\log y+1}{y^2(\log y)^2}<0\qquad(y>1).
\]
Thus $G$ is concave, and its value at every point of a closed interval is at
least the smaller of its two endpoint values.  The endpoint checks use the
following strict rational bounds:
\begin{center}
\begin{tabular}{@{}cp{.72\textwidth}@{}}
\toprule
endpoint & certified inequalities \\
\midrule
$7560$ & $\euler>577/1000$, $\log\log\log7560>3917/5000$,
  $\log\A(5040)<269/200$, and
  $\eps_0\log(3/2)<203/13500$ \\
$720720$ & $\euler>72/125$, $\log\log\log720720>191/200$,
  $\log\A(5040)<269/200$, and $\eps_0\log143<5/27$ \\
\bottomrule
\end{tabular}
\end{center}
Here $\A(5040)=403/105$, and the arguments of the last two logarithms arise
from $7560/5040=3/2$ and $720720/5040=143$.  Substitution makes each endpoint
gap a strict rational inequality, so concavity proves positivity throughout
the interval.
\end{proof}

For the much larger interval above $720720$, a compact rational certificate
replaces integer-by-integer factorization.  A certificate row consists of a
log-height interval $[L,U]$, a rational slope $\lambda$, prime-layer cutoffs,
an upper bound for the corresponding product, a lower height bound, and
lower bounds for $\log L$ and $\log U$.  The layer inequalities certify which
prime-power events can have positive reduced weight.  The product inequality
bounds the maximum possible value of
$\log\A(n)-\lambda\log n$.  Finally, concavity reduces the Robin comparison
on $[L,U]$ to two endpoint inequalities.

More precisely, for a cutoff vector $c=(c_1,\ldots,c_d)$ define the layer box
\[
 \mathcal B(c)=\{(p,j):p\ \hbox{prime},\ 1\le j\le d,\ p\le c_j\}.
\]
The reduced weight of an event is
\begin{equation}\label{eq:row-reduced-weight}
 w_{p,j}(\lambda)=g_{p,j}-\lambda\log p
 =\log p\,(\tau_{p,j}-\lambda).
\end{equation}
To state the integer boundary checks without pretending that a cutoff is
prime, extend this expression to every real base $x>1$ by
\[
 \widetilde w_j(x;\lambda)=
 \log\!\left(1+\frac{x-1}{x(x^j-1)}\right)-\lambda\log x.
\]
For a prime $p$, this is exactly $w_{p,j}(\lambda)$.  For every layer $j$,
the row proves
\[
 \widetilde w_j(c_j;\lambda)\ge0,
 \qquad \widetilde w_j(c_j+1;\lambda)\le0,
\]
using the rational logarithm bounds below.  The gain decreases with the base,
so these two boundary signs classify every prime in that layer.  The extra
check $\widetilde w_{d+1}(2;\lambda)\le0$, together with monotonicity in the
layer and base, excludes every event above the stored vector.  Thus events inside
$\mathcal B(c)$ have nonnegative weight and all events outside have
nonpositive weight.  Any integer's actual exponent packet is consequently
bounded above by the sum over this finite box; this is the precise sense in
which the box is a valid worst case.

Appendix~\ref{app:certificate} prints the complete first row, the scaled
logarithm checker, and the two endpoint predicates.  Appendix
\ref{app:cover-inventory} lists all 36 intervals and their main structural
parameters.  Keeping those arithmetic audits separate here leaves the main
argument focused on the row schema and its soundness theorem.

The logarithm certificates used here are elementary and explicit.  For a
rational $q\ge1$, put $z=(q-1)/(q+1)$.  Then $0\le z<1$ and, for $N\ge1$,
\begin{align}
 L_N(q)&=2\sum_{i=0}^{N-1}\frac{z^{2i+1}}{2i+1},\notag\\
 U_N(q)&=L_N(q)+\frac{2z^{2N+1}}{1-z^2},\label{eq:log-series}\\
 L_N(q)&\le\log q\le U_N(q).\notag
\end{align}
The upper remainder is deliberately coarse: it follows by deleting the
denominators $2i+1$ from the positive tail of the atanh series.  Consequently
both bounds in \eqref{eq:log-series} are rational and their soundness is proved
once for arbitrary $q$ and $N$.  Every row in this proof uses $N=20$.

\paragraph{Certificate discovery versus certificate checking.}
A practical way to find a row is to choose rational $L,U,\lambda$ and then,
for layer $j$, choose an integer cutoff $c_j$.  The checker compares the
threshold formula at $c_j$ and $c_j+1$: the former must have nonnegative
reduced weight and the latter nonpositive reduced weight.  Monotonicity in
the base then proves the signs for every prime in that layer.  A final check
at base $2$ proves that no still higher layer can contribute positively.
From the resulting finite box one computes the exact integer products
\[
A=\prod(p^{j+1}-1),\qquad
B=\prod p(p^j-1),\qquad
Q=\prod p.
\]
All three products are over $\mathcal B(c)$.  Their logarithms give exactly
\[
\sum_{(p,j)\in\mathcal B(c)}w_{p,j}(\lambda)
 =\log A-\log B-\lambda\log Q.
\]
For example, Row~00 has $c=[17,5,3,2,2]$.  Its layers contain respectively
the prime bases
\[
 \{2,3,5,7,11,13,17\},\quad \{2,3,5\},\quad
 \{2,3\},\quad \{2\},\quad \{2\}.
\]
Direct multiplication gives
\begin{align*}
 A&=8\,490\,305\,136\,604\,741\,632\,000,\qquad
 B=1\,651\,300\,226\,181\,365\,760\,000,\\
 Q&=367\,567\,200.
\end{align*}
The kernel then verifies, among the remaining row predicates,
\[
 A\le\frac{5141587823941}{10^{12}}B,
 \qquad
 2^{28}\frac{11486475}{8388608}\le Q.
\]
Rational bounds of the forms $A\le(\text{gain upper})B$ and
$2^e m\le Q$, together with \eqref{eq:log-series}, produce the objective and
height bounds.  If either endpoint inequality fails, one changes the slope
or shortens the interval; accepted endpoints are then chained exactly.
This search is only certificate discovery.  It need not be trusted or even
reproduced: the stored row contains all of its outputs, the general soundness
lemma proves what successful checks imply, and Lean's kernel evaluates every
row predicate from those visible rational and integer values.

The discovery process was iterative, not an optimization computation.  An
earlier checked cover contained 39 rows.  The cancellation-aware coefficient
revision in \commitlink{208b44750090b7b0e2254e3d511b4d119052cd03}{208b447}
lowered the analytic handoff from $100000$ to $74500$ and removed the last
three rows, leaving the present 36-row witness.  Rejected candidate rows were
not retained as a stable trace, so no numerical count of trial proposals is
claimed.  Nor is 36 claimed minimal: it is an upper bound supplied by one
accepted cover, and there is no nontrivial lower bound or optimality theorem
for this certificate format.  Different slopes, interval widths, or
logarithm approximants could change the row count without changing the
soundness argument.

The historical discovery runs were not instrumented per candidate, so the
paper does not report invented search times or a trial count.  Appendix
\ref{app:cover-inventory} instead reports a reproducible quantity: the elapsed
time for each accepted row module in one identified clean CI build.  Those
times measure elaboration and kernel checking on that runner, not intrinsic
certificate complexity; they vary with hardware, toolchain, and cache state.

The height data certify that the relevant prime-base product is at least
$2^e m$; the checker turns this into a lower logarithmic height with a
20-term rational log series.  The gain bound is likewise converted into an
upper logarithmic objective.  Thus this row checks the two literal rational
inequalities displayed in Lemma~\ref{lem:row-sound}, not floating-point
approximations.  Its complete definition and kernel proof are
\codefile{Robin1984/Finite/Certificates/Row00.lean}.  The other rows have the
same schema; the full list and gap-free cover proof are
\codefile{Robin1984/Finite/Certificates/AllRows.lean}, with all individual
records in \codedir{Robin1984/Finite/Certificates}.

\begin{lemma}[Soundness of a finite row]\label{lem:row-sound}
Suppose a row $[L,U]$ satisfies its layer-sign, product, height, and endpoint
inequalities.  If $n>5040$ and $L\le\log n\le U$, then $\M(n)<0$.
\end{lemma}

\begin{proof}
The sign checks and \eqref{eq:reduced-weight} bound the actual factorization
packet by the finite threshold box.  The exact product bounds convert that
box into an upper bound $B$ for
$\log\A(n)-\lambda\log n$.  The two endpoint checks state
\[
  B<\euler+\log\log L-\lambda L,
  \qquad
  B<\euler+\log\log U-\lambda U.
\]
Concavity of $y\mapsto\log\log y-\lambda y$ gives the same strict inequality
at $y=\log n$.  Rearrangement is $\M(n)<0$.
\end{proof}

\begin{proposition}[Certified finite cover]\label{prop:finite-cover}
Robin's inequality holds whenever
\begin{equation}\label{eq:finite-log-range}
  720720\le n,
  \qquad
  \log n<74500.
\end{equation}
\end{proposition}

\begin{proof}
Thirty-six rational rows form a gap-free chain from $336/25$ to
$37283397387/500000$, an endpoint strictly larger than $74500$.  Since
$336/25<\log720720$, they cover \eqref{eq:finite-log-range}.  Every row
satisfies the hypotheses of Lemma~\ref{lem:row-sound}.  In the final retained
row the first-layer product includes every prime at most $70849$; it is split
into 105 contiguous blocks, whose exact products concatenate to the required
prefix.  The blocks cover the half-open base range $[0,70850)$ without a gap,
and each block proves its numerator, denominator, and prime-base product.
Concatenation therefore reconstructs the same first-layer factors in $A,B,Q$
that a monolithic product would contain; the split changes evaluation cost,
not the mathematical packet.
All comparisons are integer or rational inequalities.  Thus the list of rows
and their checks constitute a finite proof of the proposition.
\end{proof}

Combining Propositions~\ref{prop:small-factor},
\ref{prop:startup-tangent}, and~\ref{prop:finite-cover} gives the following.

\begin{theorem}[Finite completion]\label{thm:finite}
If $5040<n$ and $\log n<74500$, then Robin's inequality holds.
\end{theorem}

In particular, Proposition~\ref{prop:ca-reduction} is now unconditional.

\section{The implication under the Riemann hypothesis}
\label{sec:rh-forward}

Let
\[
  \vartheta(x)=\sum_{p\le x}\log p,
  \qquad
  \psi(x)=\sum_{p^k\le x}\log p.
\]
The analytic estimate is most transparent after setting $H=\log n$ and
\begin{equation}\label{eq:natural-scale}
  S(H)=H^{-1/2}(\log H)^{-1}.
\end{equation}

The section is easiest read backward from its last line.  We seek
$\M(n)<0$.  Proposition~\ref{prop:budget} bounds this margin by an analytic
debit $\E(H)$ minus an arithmetic credit $\Pcost(H)$.  The weighted explicit
formula controls the debit; a complete block of prime squares supplies the
credit; and Proposition~\ref{prop:coefficient-gap} proves that, once
$H\ge74500$, the credit is at least $(113/50)S(H)$ while the debit is
strictly smaller.  The formulas below establish these three statements but
do not add another logical layer.

For a first reading, Lemma~\ref{lem:weighted-ef} may be used as the analytic
input and its kernel-estimate proof skipped; the narrative resumes at the
definition of $\Pcost(H)$ and $\E(H)$.  The intervening integration-by-parts
calculation is retained because it explains the exact constant $c_0$, the
asymmetric error bound, and the otherwise delicate endpoint $m=1$.

We first record the weighted explicit-formula input.  For $m\ge1$, define
\begin{equation}\label{eq:weight}
  w_m(t)=t^{-m-1}\frac{m\log t+1}{(\log t)^2},
  \qquad
  I_m(x)=\int_x^{\infty}(\psi(t)-t)w_m(t)\,dt.
\end{equation}
For a nontrivial zero $\rho$ of $\zeta$, counted with multiplicity, put
\[
  K_m(\rho,x)=\int_x^{\infty}
    t^{\rho-m-1}\frac{m\log t+1}{(\log t)^2}\,dt.
\]

\begin{lemma}[Robin's weighted explicit formula]\label{lem:weighted-ef}
Assume $\RH$.  For $m\ge1$ and $x\ge2$,
\begin{equation}\label{eq:weighted-ef}
  I_m(x)=-\sum_{\rho}\frac{K_m(\rho,x)}{\rho}-T_m(x),
\end{equation}
where the pole, gamma-factor, and trivial-zero correction is the explicit
nonnegative integral
\begin{equation}\label{eq:trivial-correction}
 T_m(x)=\int_x^\infty w_m(t)
 \left(\log(2\pi)+\frac12\log(1-t^{-2})\right)\,dt.
\end{equation}
For $t\ge2$ the parenthesis lies between $0$ and $\log(2\pi)$, and hence
\begin{equation}\label{eq:trivial-correction-bound}
 0\le T_m(x)\le
 \frac{\log(2\pi)x^{-m}}{\log x}.
\end{equation}
If
\[
  c_0=\euler+2-\log(4\pi),
\]
then
\begin{equation}\label{eq:weighted-bound}
 -B_m(x)-\frac{\log(2\pi)x^{-m}}{\log x}
 \le I_m(x)\le B_m(x),
\end{equation}
where
\begin{equation}\label{eq:Bmx}
 B_m(x)=c_0x^{1/2-m}(\log x)^{-1}
 \left(m+(\log x)^{-1}
 +\frac{4}{(2m-1)(\log x)^2}\right).
\end{equation}
\end{lemma}

\begin{proof}
Apply the explicit formula to the continuous cutoff whose Mellin transform
is $K_m(s,x)/s$.  The arithmetic side is the von Mangoldt sum associated
with $I_m(x)$; moving the vertical line separates the pole at $1$, the
nontrivial zeros, and the negative even zeros.  Every interchange is
justified by an explicit summable norm majorant, so the zero sum retains
analytic multiplicity rather than choosing an enumeration without proof.

The kernel calculation is not a black box.  Since
\[
 -\frac{d}{dt}\frac{t^{-m}}{\log t}
 =t^{-m-1}\frac{m\log t+1}{(\log t)^2},
\]
one integration by parts gives
\begin{equation}\label{eq:kernel-first-parts}
 \frac{K_m(\rho,x)}{\rho}
 =\frac{x^{\rho-m}}{\rho\log x}
  \int_x^\infty\frac{t^{\rho-m-1}}{\log t}\,dt.
\end{equation}
A second integration by parts, with
$d(t^{\rho-m})=(\rho-m)t^{\rho-m-1}\,dt$, separates the leading term
\[
 \frac{m}{\rho(m-\rho)}
 \frac{x^{\rho-m}}{\log x}
\]
from a remainder containing $(\log t)^{-2}$.  Bounding that remainder once
more by integration against $(\log t)^{-3}$ gives, under
$\RH$ and hence $\Re\rho=1/2$, the pointwise majorant
\[
 \left|\frac{K_m(\rho,x)}{\rho}\right|
 \le \frac1{|\rho|^2}\,
 x^{1/2-m}(\log x)^{-1}
 \left(m+(\log x)^{-1}
       +\frac4{(2m-1)(\log x)^2}\right).
\]
Hadamard factorization of the completed zeta function supplies the exact
multiplicity-counted identity
\[
  \sum_{\rho}|\rho|^{-2}=\euler+2-\log(4\pi)=c_0.
\]
Summing the majorant gives $B_m(x)$.  The combined pole/trivial-zero factor
is exactly \eqref{eq:trivial-correction}; its integrand bound and
$\int_x^\infty w_m(t)\,dt=x^{-m}/\log x$ prove
\eqref{eq:trivial-correction-bound}, explaining the asymmetry in
\eqref{eq:weighted-bound}.

The contour identity initially covers $m\ge2$.  For the endpoint let
\[
 q(t)=\frac{t(\log t+1)}{2\log t+1},\qquad
 q'(t)=\frac{\log t(2\log t+3)}{(2\log t+1)^2}.
\]
Exact integration by parts gives
\[
 I_1(x)=q(x)I_2(x)+\int_x^\infty q'(t)I_2(t)\,dt.
\]
The same identity is proved separately for the zero series and for
$T_m$, with absolute-integrability bounds licensing each interchange.
Substitution of the $m=2$ formula therefore yields \eqref{eq:weighted-ef}
at $m=1$, rather than obtaining it by a limiting convention.  This is
Robin's Lemma~2 with no omitted endpoint.
\end{proof}

Costa Pereira's estimates for $\psi-\vartheta$, with the corrected constants
from~\cite{CostaPereira1987}, control the prime-power part separated from
\eqref{eq:weighted-ef}.  Abel summation supplies the complementary lower
bound for a complete block of prime squares.

Define the complete prime cost
\begin{equation}\label{eq:prime-cost}
  \Pcost(H)=\sum_{p\le H/2}
  \min\!\left\{\frac{\log p}{H\log H},\frac1{p^2}\right\}.
\end{equation}
Also define
\begin{align}
  \E(H)={}&(2+c_0)\frac{H^{-1/2}}{\log H}
  +(c_0-2)\frac{H^{-1/2}}{(\log H)^2}\notag\\
  &+(8+4c_0)\frac{H^{-1/2}}{(\log H)^3}
  +2\frac{H^{-2/3}}{\log H}
  +\log(2\pi)\frac{H^{-1}}{\log H}.
  \label{eq:mertens-scalar}
\end{align}
Thus $\E(H)$ is a uniform upper bound assembled from the weighted
prime-function errors, whereas $\Pcost(H)$ is a nonnegative deduction forced
by completing the exponent sum to all small primes.  Their names record
their roles in \eqref{eq:budget}; neither is a new number-theoretic function
outside this proof.

\begin{proposition}[Complete exponent budget]\label{prop:budget}
Assume $\RH$.  If $n\ge1$ and $H=\log n\ge20000$, then
\begin{equation}\label{eq:budget}
  \M(n)\le \E(H)-\Pcost(H).
\end{equation}
\end{proposition}

\begin{proof}
For a prime $p$ occurring in $n$ with exponent $a$, the telescoping identity
from Lemma~\ref{lem:event-decomp} and $\log y\le y-1$ give
\begin{equation}\label{eq:prime-gain}
 \log\frac{1-p^{-(a+1)}}{1-p^{-1}}
 \le \frac1p-p^{-(a+1)}-m_p,
 \qquad
 m_p=\log(1-p^{-1})+p^{-1}\le0.
\end{equation}
Subtract $a\log p/(H\log H)$ from each prime contribution.  Completing the
sum to all primes $p\le H$ does not lose the omitted primes: if $a=0$ and
$p\le H/2$, the benchmark contribution pays at least the minimum in
\eqref{eq:prime-cost}.  Summation therefore yields a prime-reciprocal
expression minus $\Pcost(H)$.

Partial summation writes the prime-reciprocal expression as
$\log\log H$ plus the Meissel--Mertens constant and a weighted
$\vartheta$-error.
Replacing $\vartheta$ by $\psi$ separates the nonnegative prime-power tail.
Lemma~\ref{lem:weighted-ef} and the Costa Pereira bounds show that the
remaining weighted error is at most \eqref{eq:mertens-scalar}.  The Mertens
constants cancel with the sum of the $m_p$, giving exactly
\eqref{eq:budget}.
\end{proof}

The numerical heart of the large-height proof is the following strict gap.

\begin{proposition}[Explicit coefficient gap]\label{prop:coefficient-gap}
Assume $\RH$.  For $H\ge74500$,
\begin{equation}\label{eq:two-coefficients}
  \E(H)<\frac{113}{50}S(H),
  \qquad
  \Pcost(H)\ge\frac{113}{50}S(H).
\end{equation}
\end{proposition}

\begin{proof}
For the first inequality, divide \eqref{eq:mertens-scalar} by $S(H)$.
The exact identity becomes
\[
  2+c_0+(c_0-2)(\log H)^{-1}
  +(8+4c_0)(\log H)^{-2}
  +2H^{-1/6}+\log(2\pi)H^{-1/2}.
\]
Uniform rational bounds valid for every $H\ge74500$ control
$(\log H)^{-1}$, $H^{-1/6}$, $H^{-1/2}$, $c_0$, and $\log(2\pi)$.
The values actually used are
\begin{equation}\label{eq:error-rational-data}
\begin{aligned}
 \log H&\ge\frac{56}{5}, &
 (\log H)^{-1}&\le\frac5{56}, &
 (\log H)H^{-1/6}&\le\frac{87}{50},\\
 H^{-1/2}&\le\frac1{272}, &
 c_0&\le\frac1{20}, &
 \log(2\pi)&\le\frac{46}{25}.
\end{aligned}
\end{equation}
There is no numerical conflict in the bound for $c_0$:
\[
 c_0=\euler+2-\log(4\pi)
     =0.0461914179\ldots < 0.05=\frac1{20}.
\]
The proof does not depend on that decimal evaluation.  The harmonic-number
and logarithm bounds in Appendix~\ref{app:coefficients} prove
$c_0<154/3125<1/20$ rationally.  Substitution into
\eqref{eq:error-rational-data}, while retaining the favorable negative
coefficient $c_0-2$, gives the exact margin
\[
 \frac{\E(H)}{S(H)}\le\frac{3010451}{1332800}
 =\frac{113}{50}-\frac{1677}{1332800}<\frac{113}{50}.
\]

For the second inequality, take the complete prime-square block between
$\sqrt{2H}$ and $H/2$.  Abel summation gives a dominant normalized main term
tending to $2\sqrt2$, while the prime-function error, endpoint, and
prime-power terms are lower order.  The calculation retains the RH error at
the far endpoint instead of replacing the whole $\vartheta$ tail by a coarse
global bound; only after that cancellation is the remaining nonnegative
prime-power contribution discarded.  Appendix~\ref{app:coefficients} gives
the complete normalized expression and every rational substitution.  The
resulting exact lower margin is
\[
 \frac{\Pcost(H)}{S(H)}\ge
 \frac{15537277889}{6868000000}
 =\frac{113}{50}+\frac{15597889}{6868000000}>
 \frac{113}{50}.
\]
Thus the two sides of the gap are proved independently and meet with positive
rational margins.
\end{proof}

\begin{theorem}[The RH implication]\label{thm:rh-forward}
If $\RH$ holds, then Robin's inequality holds for every $n>5040$.
\end{theorem}

\begin{proof}
If $\log n<74500$, use Theorem~\ref{thm:finite}.  Otherwise combine
Propositions~\ref{prop:budget} and~\ref{prop:coefficient-gap}:
\[
  \M(n)
  <\left(\frac{113}{50}-\frac{113}{50}\right)S(\log n)
  =0.
\]
This is equivalent to Robin's inequality.
\end{proof}

\section{The Nicolas--Landau oscillation}\label{sec:oscillation}

Define Nicolas's finite product and normalized function by
\begin{equation}\label{eq:nicolas-function}
  \Pprod(x)=\prod_{p\le x}\left(1-\frac1p\right),
  \qquad
  F(x)=e^{\euler}\log\vartheta(x)\,\Pprod(x).
\end{equation}
The signed quantity used below is $\log F(x)$.

Let
\[
  k(t)=\frac{(\log t)^{-1}+(\log t)^{-2}}{t^2},
\]
and put
\begin{equation}\label{eq:KJ}
 K(x)=\int_x^{\infty}(\vartheta(t)-t)k(t)\,dt,
 \qquad
 J(x)=\int_x^{\infty}(\psi(t)-t)k(t)\,dt.
\end{equation}

For functions $g,h:\R\to\R$, write $g=\Omega_-(h)$ at infinity if there
is $c>0$ such that, for every $X$, some $x\ge X$ satisfies
$g(x)\le-c h(x)$.

Only two outputs of this section are needed later.  First,
$\log F(x)$ is bounded above by the weighted prime error $J(x)$ plus a term
of order $x^{-1}$.  Second, if $\RH$ fails, Landau's principle forces
$J(x)$ below a negative multiple of $x^{-b}$ arbitrarily far out for some
$b<1/2$.  Since $x^{-1}=o(x^{-b})$, the first output transfers the second
one to $\log F$.  Section~\ref{sec:lcm} then absorbs the remaining LCM
losses, which are only of order $P^{-1/2}\log P$.

\begin{lemma}[Finite Nicolas comparison]\label{lem:nicolas-comparison}
For all sufficiently large $x$,
\begin{equation}\label{eq:nicolas-upper}
  \log F(x)\le J(x)+\frac4x.
\end{equation}
\end{lemma}

\begin{proof}
The exact decomposition, with no asymptotic notation, is
\begin{align}
 \log F(x)={}&R_M(x)+R_\vartheta(x)+K(x),
 \label{eq:finite-nicolas-split}\\
 R_M(x)={}&
 \sum_{0<n\le\lfloor x\rfloor}
 \left(\log(1-n^{-1})+n^{-1}\right)-(M-\euler),\notag\\
 R_\vartheta(x)={}&
 \log\log\vartheta(x)-\log\log x
 -\frac{\vartheta(x)-x}{x\log x},\notag
\end{align}
where the summand is understood to be zero for nonprimes and $M$ is the
Meissel--Mertens constant in the convergent prime summand series.  The
explicit tail estimate gives $|R_M(x)|\le4/x$.  The inequality
$\log y\le y-1$, applied to the ratio of the two logarithmic heights, is
exactly $R_\vartheta(x)\le0$ after simplification.

Finally
\[
 J(x)-K(x)=\int_x^\infty
   \bigl(\psi(t)-\vartheta(t)\bigr)k(t)\,dt\ge0,
\]
because both factors are nonnegative.  Applying these three bounds to
\eqref{eq:finite-nicolas-split} gives \eqref{eq:nicolas-upper}.
\end{proof}

\begin{proposition}[Nicolas--Landau negative oscillation]
\label{prop:nicolas-landau}
If $\RH$ is false, then there is $b$ with
\begin{equation}\label{eq:b-range}
  0<b<\frac12
\end{equation}
such that
\begin{equation}\label{eq:J-omega}
  J(x)=\Omega_-(x^{-b}).
\end{equation}
\end{proposition}

\begin{proof}
Failure of $\RH$, together with the functional equation and conjugation
symmetry, gives a nontrivial zero $\rho$ of $\zeta$ with
$1/2<\Re\rho<1$.  On its horizontal line choose a zero $\rho_*$ with
maximal real part.  Choose
\[
 1-\Re\rho_*<b<\frac12.
\]

The contradiction is an instance of a simple positivity obstruction.
If the corrected tail $J(x)+x^{-b}$ were eventually positive, its Mellin
transform would be the Laplace transform, under $x=e^y$, of a positive
measure.  Positivity makes convergence propagate to every real point that
the analytic continuation can cross.  The zeta logarithmic derivative,
however, gives that continuation a pole at $s=1-\rho_*$, strictly inside
the propagated half-plane.  The remainder of the proof makes each part of
this obstruction precise.

Suppose \eqref{eq:J-omega} fails.  Then, beyond some $X$,
\begin{equation}\label{eq:positive-tail}
  J(x)+x^{-b}>0.
\end{equation}
Define the positive measure
\[
  d\mu_b(x)=\boldsymbol{1}_{(X,\infty)}(x)
       x^{-1}\bigl(J(x)+x^{-b}\bigr)\,dx
\]
and use $Y(x)=\log x$ as the random variable.  Its complex
moment-generating function is the Mellin transform
\begin{equation}\label{eq:landau-mgf}
  M_b(s)=\int_X^\infty x^{s-1}\bigl(J(x)+x^{-b}\bigr)\,dx.
\end{equation}
For $\Re s<0$ the integral is computed directly.  Integration by parts in
the $J$ term expresses it through the Mellin transform of $\psi(x)-x$ and
hence through $-\zeta'(1-s)/\zeta(1-s)$ after its pole at $1$ is removed.
The added power has the exact transform
\[
  \int_X^\infty x^{s-b-1}\,dx=\frac{X^{s-b}}{b-s}
  \qquad(\Re s<b).
\]
These formulas define an analytic continuation $H_b(s)$ wherever the zeta
factor is regular.

Here is the transform calculation used in that sentence.  The Frullani
identity
\[
 k(t)=\int_0^\infty (u+1)t^{-(u+2)}\,du
\]
and Fubini's theorem give, for $\Re s<0$,
\begin{equation}\label{eq:J-mellin-explicit}
 \int_3^\infty x^{s-1}J(x)\,dx
 =\int_0^\infty (u+1)
 \frac{\Phi_3(u+1-s)-3^s\Phi_3(u+1)}{s}\,du,
\end{equation}
where
\begin{align}
 \Phi_3(z)={}&
 \left(\frac{-\zeta'(z)}{z\zeta(z)}-\frac1{z-1}\right)\notag\\
 &-\int_1^3(\psi(t)-t)t^{-z-1}\,dt.
 \label{eq:psi-tail-continuation}
\end{align}
The apparent value at $s=0$ is filled by analytic continuation.  Absolute
integrability of the three-variable carrier $3<x<t$, $u>0$ licenses the two
exchanges of integration in \eqref{eq:J-mellin-explicit}; compact and
large-$u$ majorants then prove analyticity of the filled expression.  Formula
\eqref{eq:psi-tail-continuation} makes the singularity source explicit: the
finite startup integral is entire, while the logarithmic derivative records
the nontrivial zeros of $\zeta$.

We now separate the two roles of Landau's argument.  First, positivity of
$\mu_b$ makes the real convergence set of \eqref{eq:landau-mgf} convex.
At an interior point $a$, the derivatives of $M_b$ are its nonnegative
moments
\[
  \int (\log x)^k e^{a\log x}\,d\mu_b(x).
\]
If an analytic continuation has a power-series radius reaching $a+d$, its
Taylor series bounds the sum of these moments at $d\ge0$; monotone
convergence then proves actual integrability at $a+d$.  Iterating this strict
crossing shows that every real $\sigma<b$ lies in the interior of the true
integrability set.  Consequently $M_b$ itself is analytic throughout the
half-plane $\Re s<b$ and agrees there with $H_b$ by the identity theorem.

Second, put $s_0=1-\rho_*$.  Since
$\Re s_0=1-\Re\rho_*<b$, the preceding conclusion makes $M_b$ analytic, and
therefore locally bounded, at $s_0$.  Approach that point along
\[
  s_\delta=1-\rho_*-\delta,\qquad \delta>0.
\]
Maximality of $\Re\rho_*$ on its horizontal line says that
$\zeta(\rho_*+\delta)\ne0$, so $H_b(s_\delta)$ is regular for every
$\delta>0$.  At the endpoint, however, $1-s_0=\rho_*$ and the logarithmic
derivative has the zero pole; the norm of $H_b(s_\delta)$ tends to infinity
as $\delta\downarrow0$.  Since $M_b(s_\delta)=H_b(s_\delta)$, analytic
boundedness and divergence contradict one another.  Thus the assumed
eventual positivity is impossible, proving \eqref{eq:J-omega}.
\end{proof}

The proof just given is the strict form of the oscillation argument used by
Nicolas and Robin.  Its quantifiers are important: the excursions occur
arbitrarily far out, rather than at one selected numerical frontier.
The positive-transform step is not imported as an unnamed principle.
Its reusable formal statement proves that if a nonnegative exponential
transform converges at every real exponent below a frontier and has an
analytic continuation across the frontier, then it actually converges at
some strictly larger exponent.  The proof expands the continuation into its
power series, identifies its coefficients with nonnegative moments, and
uses monotone convergence to recover exponential integrability.  This
generic theorem is linked explicitly in Section~\ref{sec:formalization};
the subsequent Nicolas modules supply \eqref{eq:J-mellin-explicit}, the
rightmost-zero ray, and the blow-up contradiction.

\begin{corollary}\label{cor:finite-nicolas-omega}
If $\RH$ is false, there exists $0<b<1/2$ such that
\begin{equation}\label{eq:logF-omega}
  \log F(x)=\Omega_-(x^{-b}).
\end{equation}
\end{corollary}

\begin{proof}
Use Lemma~\ref{lem:nicolas-comparison} and
$x^{-1}=o(x^{-b})$, which holds because $b<1$.
\end{proof}

\section{Least common multiples and the converse}\label{sec:lcm}

Let
\[
  L_P=\operatorname{lcm}(1,2,\dots,P).
\]
The classical identity $\log L_P=\psi(P)$ connects this finite integer with
the analytic frontier.  Define the logarithmic height loss
\begin{equation}\label{eq:height-loss}
  D(P)=\log\log\psi(P)-\log\log\vartheta(P).
\end{equation}
The saturated Euler product differs from the actual abundancy of $L_P$ by
the nonnegative prime-power reserve
\begin{equation}\label{eq:tower-reserve}
  T(P)=\sum_{p\le P}
   -\log\!\left(1-p^{-a_p-1}\right),
  \qquad
  a_p=\max\{a:p^a\le P\}.
\end{equation}

\begin{lemma}[Exact LCM identity]\label{lem:lcm-identity}
For all sufficiently large $P$,
\begin{equation}\label{eq:lcm-identity}
  \M(L_P)=-\log F(P)-D(P)-T(P).
\end{equation}
\end{lemma}

\begin{proof}
The finite Euler product gives
\[
  \sum_{p\le P}-\log(1-p^{-1})
   =-\log \Pprod(P),
\]
where $\Pprod(P)$ denotes the product in \eqref{eq:nicolas-function}.  The
prime-power factorization of $L_P$ gives
\[
  -\log \Pprod(P)=\log\A(L_P)+T(P).
\]
Now expand $\log F(P)$ and use $\log L_P=\psi(P)$.  Rearrangement is exactly
\eqref{eq:lcm-identity}.
\end{proof}

\begin{lemma}[The transfer losses]\label{lem:transfer-loss}
For every $b<1/2$,
\begin{equation}\label{eq:loss-little-o}
  D(P)+T(P)=o(P^{-b}).
\end{equation}
\end{lemma}

\begin{proof}
The elementary bound
\[
  0\le\psi(P)-\vartheta(P)
  \le2\sqrt P\log P
\]
and a linear lower bound for $\vartheta(P)$ give
\[
  0\le D(P)\ll \frac{\log P}{\sqrt P}.
\]
Splitting the reserve at $\sqrt P$ and summing the geometric tails gives
$T(P)\ll P^{-1/2}$.  Since
$P^{-1/2}\log P=o(P^{-b})$ whenever $b<1/2$, the claim follows.
\end{proof}

\begin{theorem}[The converse implication]\label{thm:converse}
If Robin's inequality holds for every $n>5040$, then $\RH$ holds.
\end{theorem}

\begin{proof}
Suppose $\RH$ is false.  Choose $b$ as in
Corollary~\ref{cor:finite-nicolas-omega}.  Negative excursions of
$\log F(P)$ give positive excursions of $-\log F(P)$ of order $P^{-b}$.
By Lemma~\ref{lem:transfer-loss}, subtracting $D(P)+T(P)$ preserves such
excursions.  Lemma~\ref{lem:lcm-identity} therefore gives arbitrarily large
$P$ with $\M(L_P)>0$.  Since $L_P\ge P$, these least common multiples
eventually exceed $5040$, and $\M(L_P)>0$ is a literal failure of Robin's
inequality.  This contradicts the hypothesis.
\end{proof}

\section{Proofs of the main theorems}

\begin{proof}[Proof of Theorem~\ref{thm:main}]
The implication $\RH\Rightarrow$\eqref{eq:robin} is
Theorem~\ref{thm:rh-forward}.  The reverse implication is
Theorem~\ref{thm:converse}.
\end{proof}

\begin{proof}[Proof of Theorem~\ref{thm:ca-main}]
If $\RH$ holds, Theorem~\ref{thm:main} gives Robin's inequality for every
integer above $5040$, hence in particular for every colossally abundant one.
Conversely, assume it for every colossally abundant integer above $5040$.
The finite startup theorem and Proposition~\ref{prop:ca-reduction} extend it
to every integer above $5040$.  Theorem~\ref{thm:main} then gives $\RH$.
\end{proof}

\section{Formal verification and reproducibility}\label{sec:formalization}

\subsection{Public statements and the checked proof graph}

The preceding proof has been formalized in Lean~4.33.1 using
Mathlib~\cite{Lean4,Mathlib}.  The two public declarations have the
mathematical types
\begin{align*}
 &(\forall n\in\N,\ 5040<n\Rightarrow
   \sig(n)<e^{\euler}n\log\log n)
   \ \Longleftrightarrow\ \RH,\\
 &\RH\ \Longleftrightarrow\
   (\forall n,\eps,\ 5040<n\Rightarrow
   n\text{ is colossally abundant at }\eps\Rightarrow
   \sig(n)<e^{\euler}n\log\log n).
\end{align*}
They are named
\begin{center}
\small
\path{Robin1984.robin_inequality_iff_riemannHypothesis}\\
\path{Robin1984.riemannHypothesis_iff_colossallyAbundant_robin}.
\end{center}
The small Palomar challenge file fixes these statements using Mathlib only;
the solution file supplies the proof terms, so statement comparison is
independent of the internal proof graph.

Here the symbol $\RH$ is not an informal placeholder.  Expanding Mathlib's
definition, it is the proposition
\[
 \forall s\in\C,\quad \zeta(s)=0\ \Longrightarrow\
 \bigl(\nexists k\in\N,\ s=-2(k+1)\bigr)\ \Longrightarrow\
 s\ne1\ \Longrightarrow\ \Re s=\frac12.
\]
Thus the public biconditional uses the usual statement that every nontrivial
zero of the analytically continued zeta function lies on the critical line.
The converse proof uses the functional equation and conjugation symmetry to
turn a violation of this predicate into a zero with $1/2<\Re s<1$.
No alternate project definition of $\RH$ is introduced.  The exact reduction
from the negation of Mathlib's predicate to such a right-half zero is proved
in \codelink{Robin1984/NicolasLandau/NicolasLandau.lean}{NicolasLandau}.

The shortest route through the checked declarations is the following table.
An arrow here means use as a theorem, not merely an import.
\begin{center}
\small
\begin{tabular}{@{}p{.23\textwidth}p{.69\textwidth}@{}}
\toprule
mathematical branch & checked Lean route \\
\midrule
public boundary &
  \codelink{Challenge.lean}{Challenge} $\leftrightarrow$
  \codelink{Solution.lean}{Solution} by exact statement comparison \\
equivalence assembly &
  \codelink{Robin1984/Equivalence/Equivalence.lean}{Equivalence}
  combines the forward and converse declarations \\
forward implication &
  \codelink{Robin1984/Equivalence/LargeHeightRobin.lean}{LargeHeightRobin}
  above $74500$, and
  \codelink{Robin1984/Finite/FiniteComplete.lean}{FiniteComplete} below it \\
converse implication &
  \codelink{Robin1984/NicolasLandau/NicolasLandauRobinBridge.lean}
    {NicolasLandauRobinBridge}
  combines the negative oscillation with the LCM transfer \\
CA restriction &
  \codelink{Robin1984/ColossallyAbundant/CAReduction.lean}{CAReduction}
  combines startup completion with tangent coverage \\
\bottomrule
\end{tabular}
\end{center}
Thus the full biconditional does not emerge from one opaque analytic theorem:
the finite and large-height forward branches meet only in the final
case split, and the converse is a separate contradiction branch.

The two advertised biconditionals also have different quantifier surfaces
but no gap between them.  From $\RH$, Theorem~\ref{thm:main} immediately
specializes to every colossally abundant integer.  Conversely, the CA
hypothesis and Proposition~\ref{prop:ca-reduction} extend the inequality from
those extrema to every integer above $5040$; Theorem~\ref{thm:main} then
returns $\RH$.  Thus the second result is assembled from the all-integer
criterion and the independently proved extremal reduction; Sections
\ref{sec:finite} and \ref{sec:rh-forward} are not being mistaken for proofs
only about CA integers.

For the technically dense statements, the correspondence is declaration
level rather than merely thematic:
\begin{center}
\small
\begin{tabular}{@{}>{\raggedright\arraybackslash}p{.25\textwidth}
  >{\raggedright\arraybackslash}p{.67\textwidth}@{}}
\toprule
paper statement & checked Lean declaration and relation \\
\midrule
Lemma~\ref{lem:weighted-ef} &
 \path{robinPsiWeightedErrorIntegral_eq_zero_sum_correction_all} is
 \eqref{eq:weighted-ef} with \eqref{eq:trivial-correction};
 \path{robinPsiWeightedErrorIntegral_bounds_all} is
 \eqref{eq:weighted-bound}, both in
 \codelink{Robin1984/Equivalence/RobinLemmaTwo.lean}{RobinLemmaTwo} \\
Proposition~\ref{prop:budget} &
 \path{robin_log_abundancy_le_complete_prime_cost} in
 \codelink{Robin1984/Equivalence/AbundancyExponentBound.lean}
 {AbundancyExponentBound} is the same inequality after unfolding
 $\M$, $\E$, and $\Pcost$ \\
Proposition~\ref{prop:coefficient-gap} &
 \path{robinMertensWeightedScalar_lt_113_div_50} and
 \path{robin_minimum_prime_cost_ge_large_scalar} in the linked
 \codelink{Robin1984/Analytic/MertensExplicitCoefficient.lean}
 {Mertens coefficient} and
 \codelink{Robin1984/Analytic/PrimeCostExplicitCoefficient.lean}
 {prime-cost coefficient} modules are the two inequalities in
 \eqref{eq:two-coefficients} \\
Lemma~\ref{lem:nicolas-comparison} &
 \path{nicolasLogMertensOscillation_le_J_add_four_div} in
 \codelink{Robin1984/NicolasLandau/NicolasOscillation.lean}
 {NicolasOscillation} is \eqref{eq:nicolas-upper} \\
Landau crossing &
 \path{exists_integrableExpSet_gt_of_analyticContinuationAt} in
 \codelink{Robin1984/Mathlib/Probability/Moments/MGFAnalyticContinuation.lean}
 {MGFAnalyticContinuation} is the generic positive-transform theorem \\
Proposition~\ref{prop:nicolas-landau} &
 \path{exists_nicolasJ_omegaMinus_of_not_riemannHypothesis} in
 \codelink{Robin1984/NicolasLandau/NicolasLandauRightmostRay.lean}
 {NicolasLandauRightmostRay} has exactly the quantifiers in
 \eqref{eq:b-range}--\eqref{eq:J-omega} \\
Theorem~\ref{thm:converse} &
 \path{not_nativeRobinInequalityAll_of_not_riemannHypothesis} in
 \codelink{Robin1984/NicolasLandau/NicolasLandauRobinBridge.lean}
 {NicolasLandauRobinBridge} is its contrapositive \\
\bottomrule
\end{tabular}
\end{center}

\subsection{How the checked artifact was produced}

What the final proof says and how it was produced are different questions.
The former is answered by the declarations above and by the mathematical
proof in Sections~\ref{sec:extremal}--\ref{sec:lcm}.  The production process
was a source-guided reconstruction with the following stages.
\begin{enumerate}[leftmargin=1.7em,itemsep=.3em]
  \item The public statements were fixed first in
    \codelink{Challenge.lean}{Challenge}, directly over Mathlib definitions
    and without importing any project module.  This prevented later proof
    engineering from silently changing the target.
  \item Robin's argument and the cited inputs from Nicolas, Landau, Alaoglu--
    Erd\H{o}s, Rosser--Schoenfeld, and Costa Pereira were decomposed into named
    intermediate statements.  Source-specific results remained in their
    mathematical families, while reusable order, integration, logarithm,
    Mellin-transform, and prime-power lemmas were separated from the final
    theorem assembly.
  \item Analytic identities were proved symbolically before numerical
    constants were introduced.  Infinite sums and integrals were accompanied
    by explicit summability or integrability proofs; the weighted zero sums
    retained analytic multiplicity.
  \item The finite-range checker and its soundness theorem were proved before
    certificate data were accepted.  Candidate rows and prime blocks could be
    found by ordinary computation, but only their exact integer and rational
    outputs entered Lean, where kernel reduction checked them.
  \item The resulting source graph was reduced to the declarations needed by
    the advertised results and the public convenience facades.  Direct-import,
    provenance, placeholder, metadata, statement-boundary, and independent
    replay checks were then applied to the completed graph.
\end{enumerate}
Human direction selected the theorem statements, mathematical route, source
correspondences, and final assertions.  AI-assisted proof engineering helped
with implementation, refactoring, certificate organization, and audits.  The
production history is not itself a proof: correctness rests on the resulting
proof terms, exact certificates, and the validation pipeline described below.

\subsection{Automated research and formalization context system}

The AI assistance was not a single prompt followed by unchecked code
generation.  The project used a persistent, structured context system that
connected four kinds of state:
\begin{enumerate}[leftmargin=1.7em,itemsep=.3em]
  \item \emph{Research state} associated claims with exact source passages,
    bibliographic identities, chronology, and provenance classifications.
    Later reconstructions and formalization choices were kept distinct from
    results stated by Robin, Nicolas, Landau, and the other cited authors.
  \item \emph{Mathematical state} recorded exact definitions, quantifier
    signatures, sign conventions, normalizations, numerical domains, and the
    current hard lemma connecting an open branch to the advertised theorem.
    Numerical experiments were attached only as evidence for a matching
    domain, never promoted to hypotheses.
  \item \emph{Formal state} tracked the exact Lean declaration under attack,
    its direct and transitive dependencies, import providers, proof-state
    obligations, prior failed approaches, focused-build evidence, and the
    declarations on the final dependency path.
  \item \emph{Assurance state} tracked placeholder and axiom audits, provenance
    coverage, direct-import checks, certificate coverage, statement-boundary
    comparison, kernel replay, and release status.  Substantive mistakes were
    converted into explicit preflight rules before the same class of action
    could be repeated.
\end{enumerate}

Reusable task-specific \emph{skills} formed the control layer over that
state.  A skill is a versioned workflow package activated for a class of
work, not a learned mathematical oracle.  The packages used here covered:
\begin{itemize}[leftmargin=1.5em,itemsep=.25em]
  \item source and provenance review, including chronology and the distinction
    between direct source formalization and project-original glue;
  \item hard-theorem discipline, which fixes the live declaration, forbids
    replacing a universal parameter by an unjustified finite experiment, and
    separates proved results from certificate evidence and open obligations;
  \item numerical and certificate work, requiring the proposed inequality and
    its domain to be checked independently before Lean formalization and
    requiring every retained certificate to pass a general soundness theorem;
  \item repository engineering, including exact import providers, dependency
    closure, placeholder scans, focused builds, and staged-surface checks; and
  \item paper and artifact production, requiring source-to-declaration
    correspondence, exact-link checks, repeated PDF builds, extracted-text
    checks, and page-by-page visual inspection.
\end{itemize}
The context router activated the smallest applicable set and supplied each
skill with the same live theorem identity and evidence ledger.  Skills could
compose---for example, a numerical bound proposed during proof work could not
enter the paper until the numerical-domain rule, Lean proof rule, provenance
rule, and document-verification rule all agreed on its status.  Their durable
preflight and hard-stop conditions also carried lessons from a failed build
or incorrect claim into later sessions, instead of relying on conversational
memory alone.

The skill library also evolved during the project.  When a build failure,
mathematical contradiction, inaccurate claim, or user correction exposed a
substantive failure class, work on the affected path stopped.  The system
recorded four fields: the triggering condition, a check that must run before
the action is attempted again, the action now prohibited, and the corrected
action.  A repeated failure strengthened the entry with a mechanically tested
hard stop.  Rules that applied only to one declaration remained in the
project context; rules that generalized across formalizations were promoted
into, or used to revise, a reusable skill.  Subsequent sessions loaded the
revised package before acting, so failure-derived knowledge changed future
behavior even when the earlier conversation was no longer present.

\paragraph{Worked mathematical case: closing the large-height gap.}
One retained task packet fixed the upper-error and lower-cost Lean targets in
\codelink{Robin1984/Analytic/MertensExplicitCoefficient.lean}{Mertens coefficient}
and
\codelink{Robin1984/Analytic/PrimeCostExplicitCoefficient.lean}{prime-cost coefficient},
the universal
domain $H\ge74500$, the normalization $S(H)$, and the prohibition on adding
an assumption or changing the public theorem.  The prior checked route began
at $H=100000$.  The automated proposal that survived review had three parts:
retain the favorable sign of $(c_0-2)/\log H$ instead of taking an absolute
value; complete a block of prime squares between $\sqrt{2H}$ and $H/2$; and
carry the far-end RH error through the Abel-summation cancellation before
discarding nonnegative terms.  It then proposed separate rational subgoals
for every monotone factor.

Human acceptance was conditional on two independent checks: the normalized
upper and lower expressions had to be numerically true on the stated domain,
and Lean had to prove their symbolic monotonicity and exact rational margins.
The accepted output is not the proposal text but the declarations in
\codelink{Robin1984/Analytic/MertensExplicitCoefficient.lean}{MertensExplicitCoefficient}
and
\codelink{Robin1984/Analytic/PrimeCostExplicitCoefficient.lean}{PrimeCostExplicitCoefficient},
together with the fractions printed in Appendix~\ref{app:coefficients}.
Commit
\commitlink{208b44750090b7b0e2254e3d511b4d119052cd03}{208b447}
records the resulting $100000$-to-$74500$ handoff and the removal of three
now-unneeded certificate rows.  The numerical check was evidence; only the
Lean proofs and subsequent kernel checks were accepted into the artifact.

\paragraph{Worked failure-to-rule case: theorem and import boundaries.}
During submission hardening, a Comparator failure reported that the Challenge
and Solution theorem statements did not match even though their displayed
formulas appeared identical.  The failure class was dependence on implicit
or project-local declaration identity at the public boundary.  The corrected
rule requires \leanname{autoImplicit false}, permits only Mathlib imports in
\path{Challenge.lean}, defines the statement vocabulary directly there, and
requires exported declaration equality before release.  The public changes
are visible in commits
\commitlink{9612f5e}{9612f5e} and
\commitlink{160cd08}{160cd08}.

A second failure occurred after import narrowing: modules using
\leanname{Nat.lcmUpto} no longer imported its defining provider and failed
with an unknown-constant error.  The corresponding rule now resolves the
exact provider, requires a direct import in every consumer, and builds the
first affected consumer before any full CI run.  Commits
\commitlink{1d60111}{1d60111} and
\commitlink{0544629}{0544629} apply the correction.  The mechanical hard stop
is encoded in \codelink{scripts/check-imports.ps1}{check-imports.ps1}; the
statement hard stop is Comparator itself.  A later
\runlink{33661242178}{complete successful run} exercised the Lean build,
documentation, Comparator, Lean replay, NanoDa replay, and release gate.

These examples expose the task inputs, the accepted decomposition, the human
acceptance criterion, the retained source output, and the recurrence barrier.
They are the two documented workflow cases evaluated here, not a complete
telemetry record of every exploratory interaction.  No baseline run without
the system was performed, so the paper makes no causal claim that automation
reduced elapsed time, human effort, or the total number of errors.  The public
before-and-after commits, exact proof terms, check scripts, and CI logs make
the narrower claims about retained outputs and recurrence barriers testable.

This evolution is controlled rather than autonomous self-modification of the
proof environment.  Skill revisions are explicit, reviewable instruction
artifacts; they are checked on representative tasks and remain subordinate to
human direction and repository policy.  They may alter search order,
preflight checks, context selection, or stopping conditions, but they cannot
alter a theorem statement, convert evidence into a proof, whitelist a new
axiom, or bypass Lean's kernel and the independent replay gates.  The public
artifact therefore depends on the skills for engineering discipline, not for
logical soundness.

At each stage, the system selected only the source excerpts, definitions,
theorem types, dependency fragments, and earlier failure evidence relevant to
the live obligation.  This constrained context mattered more than raw context
volume: it reduced namespace confusion, accidental quantifier changes,
unjustified numerical specialization, duplicate lemmas, and work on
declarations outside the final proof graph.  Candidate formal steps were
tested first against the smallest relevant module; expensive full builds were
run only for coherent repository states.

Human control remained at the mathematical boundary.  The maintainer chose
the target theorems, sources, proof route, attribution, and public claims.
Automated agents did propose intermediate definitions and lemma decompositions,
search library interfaces, construct candidate proof terms, generate exact
certificate candidates, refactor dependencies, and run audits.  Human review
selected which definitions, statements, and proofs were retained.  Every
accepted output is visible as Lean source, rational data, provenance metadata,
or a validation log; no model response, research summary, numerical search,
confidence score, or context record is trusted by either advertised theorem.

The methodological claim is therefore limited to the architecture and the two
cases above.  Reproducing the checked result requires the public source and
certificates, not a replay of the exploratory interaction history: the build,
Comparator, and two kernel implementations check the retained artifact.

\subsection{Reusable formalization outputs}

Beyond the two final biconditionals, the development leaves several units
intended for reuse:
\begin{itemize}[leftmargin=1.5em,itemsep=.3em]
  \item a generic Landau positive-transform theorem in the
    \path{ProbabilityTheory} namespace, proving strict extension of
    exponential integrability from analytic continuation and nonnegative
    moments;
  \item multiplicity-aware zero-divisor and summability interfaces for
    completed zeta, together with explicit Mellin/zero-series interchange
    lemmas;
  \item exact rational logarithm bounds and a theorem-parametric finite-row
    checker that separates certificate discovery from proof;
  \item reusable bounds for prime-power layers, least-common-multiple towers,
    and normalized explicit coefficients; and
  \item provenance, dependency, Challenge/Solution, replay, and release gates
    for a large formal-mathematics artifact.
\end{itemize}
The first and other project-independent units are isolated under
\codedir{Robin1984/Mathlib}; their proposed upstream destinations are recorded
in \codelink{MATHLIB_PORTING.md}{MATHLIB\_PORTING.md}.  The result is therefore
both a checked reconstruction of Robin's theorem and a case study in turning
source mathematics, exploratory computation, and AI-assisted proof
engineering into a narrowly trusted formal artifact.

For an AI or machine-learning audience, the relevant contribution is not a
claim that a model resolved an open conjecture.  It is a controlled
human--agent workflow whose endpoint is falsifiable at declaration level:
fixed targets, explicit dependencies, exact certificates, recorded failure
classes, and independent kernel acceptance.  No new general-purpose Lean
tactic is claimed.  The reusable outputs are the mathematical interfaces,
certificate architecture, context-and-skill discipline, and validation
pipeline described above.

\subsection{Project-local mathematics and provenance}

The formalization status of the principal arguments is as follows.
\begin{itemize}[leftmargin=1.5em,itemsep=.35em]
  \item The Robin biconditional and CA restriction are project-local
    compositions in \codelink{Solution.lean}{Solution} and
    \codelink{Robin1984/Equivalence/Theorem.lean}{Equivalence.Theorem}.
  \item The weighted explicit formula, including $m=1$, is a project-local
    source formalization in
    \codelink{Robin1984/NicolasLandau/WeightedExplicitFormula.lean}{weighted formula}
    and \codelink{Robin1984/Equivalence/RobinLemmaTwo.lean}{$m=1$ endpoint}.
  \item The Nicolas--Landau negative oscillation is a project-local source
    formalization in
    \codelink{Robin1984/NicolasLandau/NicolasLandauRightmostRay.lean}{oscillation proof}.
  \item The LCM transfer and converse are project-local source formalizations
    in \codelink{Robin1984/NicolasLandau/NicolasLandauRobinBridge.lean}{LCM bridge}.
  \item Finite-row soundness and all row data are project-original
    certificate formalizations in
    \codelink{Robin1984/Finite/FiniteRowCertificate.lean}{row soundness}
    and \codelink{Robin1984/Finite/Certificates/AllRows.lean}{all rows}.
\end{itemize}
The dependency boundary can be summarized without inspecting import syntax:
\begin{center}
\small
\begin{tabular}{@{}l>{\raggedright\arraybackslash}p{.72\textwidth}@{}}
\toprule
source & mathematical role \\
\midrule
Mathlib &
  definitions of $\zeta$, $\xi$, and $\RH$; complex and real analysis,
  measure and integration theory, infinite sums, arithmetic, and elementary
  asymptotics \\
pinned PNTAnd fork &
  Hadamard-factorization and divisor infrastructure, completed-zeta facts,
  Mertens and prime-function infrastructure, and the explicit
  Rosser--Schoenfeld and Costa Pereira estimates \\
Robin1984 &
  the weighted formula in the form used here, its zero-series interchanges
  and $m=1$ endpoint, the complete exponent budget and coefficient gap, the
  Nicolas--Landau continuation and oscillation, the LCM transfer, the CA
  reduction, and every finite certificate \\
\bottomrule
\end{tabular}
\end{center}
It would therefore be incorrect to describe every dependency as proved from
scratch in repository-owned modules.  Mathlib and the pinned PNTAnd revision
supply a large body of foundational and explicit analytic number theory.
The project-local claim is narrower and checkable: the source-specific chain
from those imported results to Robin's two advertised biconditionals,
including the weighted endpoint, Nicolas--Landau continuation, LCM transfer,
CA reduction, and finite certificates, is proved here without adding a new
axiom.

In particular, Lemma~\ref{lem:weighted-ef} and
Proposition~\ref{prop:nicolas-landau} are not assumptions imported as opaque
theorems: their full proofs occur in the linked repository-owned modules.
The lower-level imported analysis is substantial, as it should be for a
theorem about zeta zeros, but no imported declaration states Robin's
biconditional or the Nicolas--Landau conclusion needed here.  The exact
distinction is recorded declaration by declaration in the provenance
metadata.

The formal development separates arithmetic identities, colossally abundant
extrema, the RH estimates, the Nicolas--Landau argument, finite certificates,
and final assembly.  Reusable order, logarithm, Mellin-transform, and
prime-power lemmas are exposed in their mathematical families.  Each owned
Lean file carries a content-reviewed provenance label distinguishing direct
source formalization, another published source, standard mathematics without
a claimed individual originator, and primarily project-authored
formalization work.  The complete reviewed ledger is distributed with the
repository as the \codelink{docs/PROVENANCE.md}{provenance guide} and its
\codelink{provenance/ledger.json}{machine-readable ledger}.

For more detailed reading, the assembled
\codelink{Robin1984/Finite/RobinTangentStartupComplete.lean}{startup transition},
the \codelink{Robin1984/Analytic/MertensExplicitCoefficient.lean}{Mertens coefficient},
the \codelink{Robin1984/Analytic/PrimeCostExplicitCoefficient.lean}{prime-cost coefficient},
and the \dirlink{Robin1984/Finite/Certificates}{finite certificate data}
correspond directly to Sections~\ref{sec:extremal}--\ref{sec:rh-forward}.
These source links follow the main branch for browsing; each release archives
the exact commit from which its paper and Lean artifacts were built.

\subsection{Exact finite certificates}

The finite proof contains no floating-point oracle.  The startup
factorizations and the 36 interval rows reduce to exact natural-number and
rational inequalities.  The large final retained prime product is split into
105 blocks only to keep individual kernel computations manageable.  The
certificate checkers are proved sound once; each data row is then accepted by
kernel reduction.

Appendix~\ref{app:certificate} quotes the checked interface and expands Row~00
into its stored record, scaled logarithm bounds, product tests, and endpoint
predicates.  The conceptual separation is simple: one ordinary theorem proves
that any successful row implies Robin's inequality throughout its interval;
another declaration asks Lean's kernel to evaluate the visible rational row.
The aggregate theorem checks that the accepted endpoints form a gap-free
cover.  The full definitions are in
\codelink{Robin1984/Finite/FiniteRowCertificate.lean}{FiniteRowCertificate}
and \codelink{Robin1984/Finite/Certificates/Row00.lean}{Row00}.

\subsection{What Comparator and the kernel replays establish}

The validation stages have different jobs.  The ordinary Lean build checks
that every project module elaborates and that Lean's kernel accepts the final
proof terms.
\href{https://github.com/leanprover/comparator/tree/68a064109f01c08f47c8edc9f51d6a2bbffaa188}
{Comparator} then protects the public theorem boundary against a different
risk: a purported solution might prove an easier proposition with the same
informal description, or redefine a symbol occurring in the target.

The trusted \codelink{Challenge.lean}{Challenge} module imports only four
narrow Mathlib modules.  It defines the divisor sum, Robin inequality, and CA
predicate directly, and states the two advertised theorems with deliberate
proof holes.  The \codelink{Solution.lean}{Solution} module imports the project
and supplies the actual proofs.  The checked
\codelink{comparator.json}{Comparator configuration} names exactly those two
theorems, permits only
\leanname{propext}, \leanname{Quot.sound}, and
\leanname{Classical.choice}, requires NanoDa, and declares no definition
holes.  In particular, the solution is not allowed to choose a convenient
meaning for $\sig$, Robin's inequality, CA, or $\RH$.

The complete semantic policy is short enough to display:
\begin{verbatim}
"theorem_names": [
  "Robin1984.robin_inequality_iff_riemannHypothesis",
  "Robin1984.riemannHypothesis_iff_colossallyAbundant_robin"
],
"definition_names": [],
"permitted_axioms": ["propext", "Quot.sound", "Classical.choice"],
"enable_nanoda": true
\end{verbatim}
The omitted JSON fields only name the modules \path{Challenge} and
\path{Solution}; there is no additional whitelist of project definitions.

Using pinned tool revisions, the Comparator job performs the following checks.
\begin{enumerate}[leftmargin=1.7em,itemsep=.3em]
  \item It builds and exports the Challenge and Solution environments
    separately under the Landrun sandbox.
  \item For each advertised theorem it compares the exported statement and
    every declaration on which that statement depends.  This is a comparison
    of Lean declarations, not of displayed text or theorem names alone.
  \item It follows the Solution theorem bodies transitively and rejects any
    axiom outside the three permitted standard principles.
  \item It replays the exported Solution environment through Lean's default
    kernel.
  \item Because \leanname{enable\_nanoda} is true, it also submits the export
    to NanoDa, an independent kernel implementation.
\end{enumerate}
Passing Comparator therefore establishes statement identity, the configured
axiom bound, and kernel acceptance through two implementations.  It does not
establish that the trusted Challenge statement is the historically intended
reading of Robin's paper, nor that the exposition is persuasive.  The explicit
Challenge definitions, the expansion of Mathlib's $\RH$ predicate above, and
ordinary mathematical review address that separate question.

\subsection{Axioms and dependency boundary}

The repository contains no project-local axiom declarations.  Outside the
two deliberate statement holes in the Palomar challenge, project proof
sources contain no \leanname{sorry}, \leanname{admit}, or
\leanname{native\_decide}.  Diagnostic commands such as
\leanname{\#print axioms} occur neither in repository-owned Lean sources nor
in the timed external import closure.  A completed kernel audit of both
advertised declarations returned exactly the standard principles
\[
  \leanname{propext},\qquad
  \leanname{Classical.choice},\qquad
  \leanname{Quot.sound}.
\]
This audit follows the proof terms of the two advertised declarations through
their actual transitive closure.  It therefore distinguishes library files
that are merely present in a dependency checkout from declarations genuinely
used by this proof: an unused theorem elsewhere cannot enter the result's
axiom set.

The ordinary \leanname{lake build} command is the full formalization build.
Its Lake-native default target reads a checked-in topological order of 257
modules, schedules one module job at a time, bounds each Lean process to one
internal task thread, and ends with \leanname{Solution}.  Thus the submitted
project itself controls the memory-intensive finite-certificate order; no
external runner or local patching step is part of the proof build.

All dependencies are pinned.  The required Lean~4.33.1 changes to
\leanname{PrimeNumberTheoremAnd} are published in the public fork
\href{\pntrepositoryurl}{\nolinkurl{https://github.com/kimihiro64/PrimeNumberTheoremAnd}}
at
\href{\pntrepositoryurl/commit/\pntrevision}{commit \texttt{f8f58c7}}, based
on
\href{https://github.com/AlexKontorovich/PrimeNumberTheoremAnd/commit/47fa48680663df41146704d02a5b092d792bd5b9}
{upstream commit \texttt{47fa486}}.  The full hashes are recorded in the Lake
manifest.  Lake reconstructs this dependency directly from public Git
history; no local patching step is required.

The fork is intended as a temporary compatibility and extension branch, not
as a permanent divergence.  At the time of this revision, the
\href{https://github.com/AlexKontorovich/PrimeNumberTheoremAnd/blob/main/lean-toolchain}
{upstream \leanname{main} toolchain file} still declares Lean~4.32.2, and
upstream \leanname{main} does not yet contain the xi-divisor source used here.
Replacing the fork therefore requires two things: an upstream revision
compatible with the project's Lean toolchain, and upstream versions or
equivalent replacements for the additional xi-divisor and Mertens interfaces.

The fork has two distinct roles.  First, it updates proof scripts and package
pins for Lean~4.33.1 and exposes named consequences of the existing Mertens
development needed by this project.  These changes repair elaboration against
the new toolchain; they do not replace an analytic theorem by an assumption
or alter Mathlib's definition of $\RH$.  Second, it adds the
multiplicity-aware interface used by the weighted zero sums.  The principal
declarations in
\pntfile{PrimeNumberTheoremAnd/Mathlib/NumberTheory/LSeries/RiemannXiDivisorZeros.lean}
  {RiemannXiDivisorZeros}
are:
\begin{itemize}[leftmargin=1.5em,itemsep=.25em]
  \item \path{RiemannXiDivisorZeroIndex} and
    \path{riemannXiDivisorZeroValue} give an index type and value map for the
    nonzero divisor of $\xi$, retaining analytic multiplicity;
  \item \path{riemannXiDivisorZeroValue_ne_zero} and
    \path{riemannXiDivisorZeroValue_eq_zero} prove that every indexed value
    is nonzero and is an actual zero of $\xi$;
  \item \path{summable_riemannXiDivisorZero_norm_inv_sq} proves summability
    of the inverse-square norms of those values;
  \item \path{completedRiemannZeta_eq_cpow_mul_Gamma_mul_riemannZeta}
    expands completed zeta under its stated nonvanishing hypotheses; and
  \item \path{neg_two_mul_logDeriv_riemannXi_zero_eq} proves the exact identity
    $-2(\xi'/\xi)(0)=\euler+2-\log(4\pi)$.
\end{itemize}
The companion
\pntfile{PrimeNumberTheoremAnd/IEANTN/RobinXiZeroSummability.lean}
  {RobinXiZeroSummability}
proves, as
\path{summable_riemannXiDivisorZero_norm_inv_rpow}, the fractional
inverse-norm summability used to interchange the zero sum and Mellin
integral.  These results are derived from the dependency's
existing divisor support, Hadamard growth and summability, completed-zeta,
and differentiability theorems.  The fork adds no axiom, proof placeholder,
or numerical oracle.  The xi-divisor files retain Matteo Cipollina's
authorship and Apache-2.0 notice; the upstream project is due to Alex
Kontorovich, Terence Tao, and its contributors.  The complete
\href{https://github.com/kimihiro64/PrimeNumberTheoremAnd/compare/47fa48680663df41146704d02a5b092d792bd5b9...f8f58c749d6cde8a641348fcd5e4702993651cd6}
{fork diff} is public.

Project-independent helper results are isolated under
\codedir{Robin1984/Mathlib} and mapped to proposed upstream Mathlib paths in
\codelink{MATHLIB_PORTING.md}{MATHLIB\_PORTING.md}.  The inventory records
which files are ready but not yet submitted.  The intended maintenance path
is to upstream those generic units separately and to propose the xi-divisor
interface and compatibility repairs to \leanname{PrimeNumberTheoremAnd};
until acceptance, the pinned public fork is the reproducible source of the
checked dependency.  When both replacement conditions are met, the project
can repin directly to upstream, but only after the full Lean build and
Comparator/NanoDa pipeline pass against that revision.

\subsection{Reproducing and independently checking the result}

The audit record is public and readable without repository membership or
GitHub authentication:
\begin{description}[leftmargin=1.7em,style=nextline,itemsep=.25em]
  \item[Source and history.]
    \url{https://github.com/kimihiro64/Robin1984}
  \item[Complete CI runs and job logs.]
    \url{https://github.com/kimihiro64/Robin1984/actions/workflows/ci.yml}
  \item[Tagged source, paper, API site, compiled Lean tree, and checksums.]
    \url{https://github.com/kimihiro64/Robin1984/releases}
  \item[Machine-readable theorem and provenance metadata.]
    \codefile{formalization.yaml}, \codefile{comparator.json}, and
    \codefile{provenance/ledger.json}
  \item[Trusted statement boundary and implementation.]
    \codefile{Challenge.lean} and \codefile{Solution.lean}
\end{description}
The string \path{kimihiro64} is simply the public GitHub owner namespace,
not an indication of a private repository.  Each release is tied to an exact
Git commit and is therefore a citable immutable snapshot even though the
\path{main} branch continues to develop.  A DOI-bearing archival mirror would
add another persistence layer, but is not required to inspect or reproduce
the present artifact.

There are three useful levels of reproduction.  First, a reader can inspect a
completed run without compiling Lean.  The
\href{\actionsurl}{public CI history} exposes the separate job results and
logs.  After every required job succeeds, the workflow publishes a
\href{\releasesurl}{GitHub release} containing the paper, an offline API site,
and the complete Linux \leanname{.lake/build} tree.  The release identifies
the exact source commit and records checksums for its assets.  The compiled
objects are a convenience for inspection and compatible incremental use; the
source and pinned revisions remain authoritative.

Second, a source-level kernel rebuild can be made from a clean release tag:
\begin{verbatim}
git clone https://github.com/kimihiro64/Robin1984.git
cd Robin1984
git checkout TAG
lake update
lake exe cache get
lake build
\end{verbatim}
Here \leanname{TAG} is the tag named on the chosen GitHub release.
\leanname{lake update} reconstructs the exact dependency revisions from the
manifest, and \leanname{lake exe cache get} downloads the matching Mathlib
cache; omitting the cache step gives a larger cold build.  The final
\leanname{lake build} follows the repository's checked-in sequential module
order and ends by building \leanname{Solution}.  On Windows,
\codelink{scripts/bootstrap.ps1}{bootstrap.ps1} installs and checks the same
project prerequisites before these commands.

The repository-level structural checks can then be run as
\begin{verbatim}
pwsh -File scripts/check-imports.ps1
pwsh -File scripts/check-provenance.ps1
ruby scripts/validate-formalization.rb
\end{verbatim}
They check narrow imports, the candidate-module boundary, per-file provenance
and documentation, and the Palomar metadata.  They are consistency checks,
not substitutes for kernel checking.

Third, on Linux with Git, Lake, Go, Rust/Cargo, Python~3, and Landrun, the
independent statement and replay pipeline is
\begin{verbatim}
./scripts/verify-comparator.sh
\end{verbatim}
The checked-in
\codelink{scripts/verify-comparator.sh}{verification script}
checks out exact reviewed revisions of Comparator, lean4export, Landrun, and
NanoDa; requires the project's Lean~4.33.1 toolchain; performs the sequential
project build; and runs the five Comparator stages listed above.  The
checked-in
\codelink{.github/workflows/ci.yml}{CI workflow} is the complete executable
recipe: it additionally validates licensing and release metadata, builds the
paper and API documentation, transfers the completed Lean artifacts to the
downstream jobs, and publishes a release only if every dependency job,
including Comparator, succeeds.

\paragraph{Observed cost and hardware.}
The hosted workflow uses GitHub's standard public
\href{https://docs.github.com/en/actions/reference/runners/github-hosted-runners}
{\path{ubuntu-latest} runner}, which at the time of this revision provides
four x64 CPU cores, 16~GB RAM, and 14~GB SSD storage.  This is a tested
sufficient environment, not a proven minimum; peak memory below 16~GB and a
smaller storage floor have not been benchmarked.  The project build is
deliberately serialized to control memory, so a clean \leanname{lake build}
should be budgeted at roughly two hours; one measured cold build took about
1~hour 42~minutes.

The independent pipeline is longer.  In the representative source-changing
\runlink{33661242178}{successful run of commit \texttt{4d0b59e}}, the Lean
job took 1~hour 16~minutes, API documentation 16~minutes 58~seconds,
Comparator and its two replays 2~hours 39~minutes, and total wall time was
3~hours 56~minutes 59~seconds because downstream jobs overlap.  In a
\runlink{33477843519}{cache-reusing successful run}, the Lean job fell to
15~minutes 36~seconds, while Comparator still took 2~hours 17~minutes.
Accordingly, an independent team should budget about four hours for a full
uncached assurance run, not merely the duration of \leanname{lake build}.
Source changes invalidate the project-object cache; paper-only changes can
reuse it.  The release's compiled tree avoids the first rebuild for a matching
environment but does not replace source-level verification.

\paragraph{Portability and maintenance.}
The authoritative artifact is the source at an immutable release tag.
The downloadable \leanname{.olean} tree is a Linux x64 convenience built for
Lean~4.33.1; compiled objects are not expected to work across Lean versions,
platforms, or incompatible dependency revisions.  A source rebuild uses
\codefile{lean-toolchain} and \codefile{lake-manifest.json}, which currently
pin Mathlib commit
\texttt{0df444a360eaa60ab8c11dca51a86af692955474} and the PNTAnd fork commit
\texttt{f8f58c749d6cde8a641348fcd5e4702993651cd6}.  The Windows bootstrap covers
the main Lean build; the full Comparator/Landrun/NanoDa recipe is presently
Linux-oriented and additionally requires the tools listed above.

Maintenance is deliberate rather than automatic.  A toolchain or dependency
repin is tested first on the affected modules and then must pass the complete
Lean, provenance, Comparator, Lean-kernel, and NanoDa gates before release.
Generic project-owned modules are tracked for upstreaming in
\codelink{MATHLIB_PORTING.md}{MATHLIB\_PORTING.md}.  As verified for this
revision, PNTAnd upstream \path{main} still declares Lean~4.32.2 and lacks the
xi-divisor interface, so the public fork remains necessary.  The fork will be
removed only when both compatibility and interface replacements exist and the
full pipeline passes after repinning.  Future maintenance may be required as
Lean evolves; immutable tagged releases preserve the exact historical inputs
rather than promising that old compiled objects will load in new toolchains.

The formalization was developed with human mathematical direction and
AI-assisted proof engineering.  Lean's kernel checks that the linked proof
terms inhabit the exact public types above, using only the reported logical
principles.  It does not by itself certify that those types are the intended
reading of Robin's article, that the reconstruction is historically faithful,
that the result is novel, or that the exposition is persuasive.  Those are
mathematical and scholarly review obligations.  The explicit expansion of
Mathlib's $\RH$ predicate, the source-to-module table, and the provenance
ledger make that comparison inspectable rather than implicit.  The present
status is self-assessed; no independent expert review is claimed.

\appendix

\section{Exact rational coefficient audit}\label{app:coefficients}

This appendix collects the arithmetic intentionally omitted from the main
narrative of Proposition~\ref{prop:coefficient-gap}.  First, the formal bound
for the zero constant uses
\[
 \euler<H_{256}-\log256,\qquad
 H_{256}\le\frac{15311}{2500},\qquad
 \log2\ge\frac{69314}{100000},\qquad
 \log(4\pi)\ge\frac{253}{100}.
\]
Consequently
\[
 c_0<
 \frac{15311}{2500}-8\frac{69314}{100000}
 +2-\frac{253}{100}
 =\frac{154}{3125}<\frac1{20}.
\]
This is the rational proof behind the decimal orientation
$c_0=0.0461914179\ldots$.

For the error bound, put $q=(\log H)^{-1}$.  The
$H^{-1/6}$ estimate in \eqref{eq:error-rational-data} turns
$2H^{-1/6}$ into at most $(87/25)q$.  Combining this with the favorable
term $-2q$ and applying the remaining bounds gives
\begin{align*}
 \frac{\E(H)}{S(H)}
 &\le
 2+\frac{37}{25}\frac5{56}
 +8\left(\frac5{56}\right)^2\\
 &\quad+\frac1{20}
   \left(1+\frac5{56}+4\left(\frac5{56}\right)^2\right)
 +\frac{46}{25}\frac1{272}\\
 &=\frac{3010451}{1332800}
 =\frac{113}{50}-\frac{1677}{1332800}
 <\frac{113}{50}.
\end{align*}

For the complete prime cost, Abel summation over the prime-square block
$\sqrt{2H}<p\le H/2$ gives the following normalized lower expression:
\begin{align*}
 &2\sqrt2\frac{\log H}{\log H+\log2}
    \left(1-\frac1{\log H+\log2}\right)\\
 &\quad-\frac94\,2^{1/4}H^{-1/4}
    \frac{\log H}{\log H+\log2}
 -\log(2\pi)H^{-1/2}
    \frac{\log H}{\log H+\log2}\\
 &\quad-4H^{-1/2}\frac{\log H}{\log H-\log2}\\
 &\quad-2\sqrt2\,c_0H^{-1}\frac{\log H}{\log H-\log2}
    \left(2+\frac1{\log H-\log2}
      +\frac4{3(\log H-\log2)^2}\right).
\end{align*}
The last line is the far-end RH term retained through the cancellation.
The rational bounds used for $H\ge74500$ are
\begin{align}\label{eq:cost-rational-data}
 \sqrt2&\ge\frac{7071}{5000}, & 2^{1/4}&\le\frac{119}{100},\notag\\
 \frac{\log H}{\log H+\log2}&\ge\frac{16}{17}, &
 \frac1{\log H+\log2}&\le\frac{17}{202},\notag\\
 H^{-1/4}\frac{\log H}{\log H+\log2}&\le\frac{573}{10000}, &
 H^{-1/2}\frac{\log H}{\log H+\log2}&\le\frac7{2000},\notag\\
 H^{-1/2}\frac{\log H}{\log H-\log2}&\le\frac{197}{50000}, &
 \frac{\log H}{\log H-\log2}&\le\frac{107}{100}.
\end{align}
The normalized far-end term is at most $1/10000$, so the final exact
substitution is
\begin{align*}
 &2\frac{7071}{5000}\frac{16}{17}
    \left(1-\frac{17}{202}\right)
 -\frac94\frac{119}{100}\frac{573}{10000}
 -\frac{46}{25}\frac7{2000}\\
 &\hspace{2em}-4\frac{197}{50000}-\frac1{10000}\\
 &\qquad=
 \frac{15537277889}{6868000000}
 =\frac{113}{50}+\frac{15597889}{6868000000}
 >\frac{113}{50}.
\end{align*}

\section{A finite certificate in full}\label{app:certificate}

The first of the 36 accepted rows stores the following exact record:
\begin{center}
\small
\begin{tabular}{@{}ll@{}}
\toprule
field & stored rational value \\
\midrule
$[L,U]$ & $[336/25,\ 22117363/1000000]$ \\
$\lambda$ & $9771914406623/500000000000000$ \\
layer cutoffs & $[17,5,3,2,2]$ \\
gain upper bound & $5141587823941/1000000000000$ \\
height data $(e,m)$ & $(28,\ 11486475/8388608)$ \\
lower bounds for $\log L,\log U$
  & $1299117667/500000000,\ 619272591/200000000$ \\
\bottomrule
\end{tabular}
\end{center}

The essential checked interface, with one proof body elided, is
\begin{verbatim}
theorem RobinFiniteRowChecks.sound {r : RobinFiniteRow}
    (h : RobinFiniteRowChecks r) {n : Nat} (hn : 5040 < n)
    (hLo : (r.lo : Real) <= Real.log (n : Real))
    (hHi : Real.log (n : Real) <= (r.hi : Real)) :
    Robin1984.Core.NativeRobinInequality n := by ...

theorem robinFiniteRow00_checks :
    RobinFiniteRowChecks robinFiniteRow00 := by
  decide +kernel
\end{verbatim}
Let $g=5141587823941/10^{12}$ be the stored gain bound and
$k(q)=\lfloor\log_2q\rfloor$.  In terms of the functions in
\eqref{eq:log-series}, the scaled upper checker is
\[
 U^{\rm sc}_{20}(q)=
 k(q)\frac{693147180559945310}{10^{18}}
 +U_{20}\!\left(\frac{q}{2^{k(q)}}\right),
\]
and the lower checker uses
$693147180559945309/10^{18}$ and $L_{20}$.  With the stored
$\lambda$, exponent $e=28$, and mantissa $m=11486475/8388608$, the
objective supplied to both endpoints is
\[
 B_0=U^{\rm sc}_{20}(g)
 -\lambda\left(
 28\frac{693147180559945309}{10^{18}}+L_{20}(m)\right).
\]
The kernel decision verifies the integer product inequalities
$A\le gB$ and $2^{28}m\le Q$, every adjacent layer-sign comparison, and
\[
 B_0<
 \frac{57721}{100000}+L^{\rm sc}_{20}(\ell_L)-\lambda L,
 \qquad
 B_0<
 \frac{57721}{100000}+L^{\rm sc}_{20}(\ell_U)-\lambda U.
\]
The general logarithm and Euler-constant theorems then convert these
rational facts into the real endpoint inequalities used by
Lemma~\ref{lem:row-sound}.  The record and its kernel proof are
\codelink{Robin1984/Finite/Certificates/Row00.lean}{Row00}; the reusable
soundness theorem is
\codelink{Robin1984/Finite/FiniteRowCertificate.lean}{FiniteRowCertificate}.

\clearpage
\section{Complete finite-cover inventory}\label{app:cover-inventory}

Table~\ref{tab:row-inventory} lists every accepted row.  The columns $L,U$
and $\lambda$ are the exact stored rationals.  The entry $c_1/d$ gives the
first-layer cutoff and the number of stored layers; ''blocks'' is the number
of separately checked first-layer prime-product blocks, with zero denoting the
direct monolithic calculation used by Row~00.  The final column is the elapsed
module time reported by the clean
\runlink{33661242178}{CI build of commit \texttt{4d0b59e}}.  These are
observed elaboration and kernel-checking times on that runner, not discovery
times or hardware-independent complexity bounds.

\begingroup
\footnotesize
\setlength{\tabcolsep}{1.5pt}
\renewcommand{\arraystretch}{1.12}
\begin{longtable}{@{}rccc rrr@{}}
\caption{The complete 36-row certificate cover.}\label{tab:row-inventory}\\
\toprule
row & $L$ & $U$ & $\lambda$ & $c_1/d$ & blocks & seconds \\
\midrule
\endfirsthead
\toprule
row & $L$ & $U$ & $\lambda$ & $c_1/d$ & blocks & seconds \\
\midrule
\endhead
00 & $\tfrac{336}{25}$ & $\tfrac{22117363}{1000000}$ & $\tfrac{9771914406623}{500000000000000}$ & 17/5 & 0 & 2.3 \\
01 & $\tfrac{22117363}{1000000}$ & $\tfrac{35788699}{1000000}$ & $\tfrac{2565510579241}{250000000000000}$ & 28/6 & 2 & 2.3 \\
02 & $\tfrac{35788699}{1000000}$ & $\tfrac{29149999}{500000}$ & $\tfrac{5519615043151}{1000000000000000}$ & 46/7 & 3 & 2.4 \\
03 & $\tfrac{29149999}{500000}$ & $\tfrac{47544747}{500000}$ & $\tfrac{3004424868557}{1000000000000000}$ & 76/7 & 4 & 2.4 \\
04 & $\tfrac{47544747}{500000}$ & $\tfrac{146801491}{1000000}$ & $\tfrac{1724212981141}{1000000000000000}$ & 120/8 & 5 & 2.5 \\
05 & $\tfrac{146801491}{1000000}$ & $\tfrac{111686563}{500000}$ & $\tfrac{8278925801}{8000000000000}$ & 184/9 & 6 & 2.5 \\
06 & $\tfrac{111686563}{500000}$ & $\tfrac{163594351}{500000}$ & $\tfrac{80829194299}{125000000000000}$ & 274/10 & 7 & 2.6 \\
07 & $\tfrac{163594351}{500000}$ & $\tfrac{464815043}{1000000}$ & $\tfrac{422180858753}{1000000000000000}$ & 395/10 & 8 & 2.6 \\
08 & $\tfrac{464815043}{1000000}$ & $\tfrac{130165559}{200000}$ & $\tfrac{70868018799}{250000000000000}$ & 557/11 & 9 & 2.7 \\
09 & $\tfrac{130165559}{200000}$ & $\tfrac{893036263}{1000000}$ & $\tfrac{38967393937}{200000000000000}$ & 771/11 & 10 & 2.7 \\
10 & $\tfrac{893036263}{1000000}$ & $\tfrac{240770673}{200000}$ & $\tfrac{137136588803}{1000000000000000}$ & 1048/12 & 12 & 2.8 \\
11 & $\tfrac{240770673}{200000}$ & $\tfrac{63176719}{40000}$ & $\tfrac{99275427189}{1000000000000000}$ & 1391/12 & 13 & 2.8 \\
12 & $\tfrac{63176719}{40000}$ & $\tfrac{2035337871}{1000000}$ & $\tfrac{73775311097}{1000000000000000}$ & 1806/13 & 14 & 2.9 \\
13 & $\tfrac{2035337871}{1000000}$ & $\tfrac{518900301}{200000}$ & $\tfrac{11152011487}{200000000000000}$ & 2314/13 & 16 & 2.9 \\
14 & $\tfrac{518900301}{200000}$ & $\tfrac{815385537}{250000}$ & $\tfrac{42786768497}{1000000000000000}$ & 2927/14 & 17 & 3.1 \\
15 & $\tfrac{815385537}{250000}$ & $\tfrac{406392661}{100000}$ & $\tfrac{6654185407}{200000000000000}$ & 3662/14 & 19 & 3.1 \\
16 & $\tfrac{406392661}{100000}$ & $\tfrac{1248690297}{250000}$ & $\tfrac{2622639723}{100000000000000}$ & 4528/14 & 21 & 3.1 \\
17 & $\tfrac{1248690297}{250000}$ & $\tfrac{3045165617}{500000}$ & $\tfrac{20930175389}{1000000000000000}$ & 5542/15 & 23 & 3.3 \\
18 & $\tfrac{3045165617}{500000}$ & $\tfrac{458996953}{62500}$ & $\tfrac{16893531041}{1000000000000000}$ & 6716/15 & 25 & 3.4 \\
19 & $\tfrac{458996953}{62500}$ & $\tfrac{4387795691}{500000}$ & $\tfrac{13794103247}{1000000000000000}$ & 8059/15 & 27 & 3.5 \\
20 & $\tfrac{4387795691}{500000}$ & $\tfrac{10433152137}{1000000}$ & $\tfrac{11354367573}{1000000000000000}$ & 9603/15 & 30 & 3.6 \\
21 & $\tfrac{10433152137}{1000000}$ & $\tfrac{12279993269}{1000000}$ & $\tfrac{2357543377}{250000000000000}$ & 11356/16 & 33 & 3.8 \\
22 & $\tfrac{12279993269}{1000000}$ & $\tfrac{179906601}{12500}$ & $\tfrac{1973616211}{250000000000000}$ & 13335/16 & 36 & 4.0 \\
23 & $\tfrac{179906601}{12500}$ & $\tfrac{16784073643}{1000000}$ & $\tfrac{6644793849}{1000000000000000}$ & 15587/16 & 39 & 4.3 \\
24 & $\tfrac{16784073643}{1000000}$ & $\tfrac{9718846411}{500000}$ & $\tfrac{1126354703}{200000000000000}$ & 18110/16 & 42 & 4.6 \\
25 & $\tfrac{9718846411}{500000}$ & $\tfrac{22433188173}{1000000}$ & $\tfrac{1200244709}{250000000000000}$ & 20934/17 & 46 & 4.9 \\
26 & $\tfrac{22433188173}{1000000}$ & $\tfrac{1287242339}{50000}$ & $\tfrac{2057220229}{500000000000000}$ & 24088/17 & 50 & 5.4 \\
27 & $\tfrac{1287242339}{50000}$ & $\tfrac{7349340853}{250000}$ & $\tfrac{3547339421}{1000000000000000}$ & 27570/17 & 54 & 5.8 \\
28 & $\tfrac{7349340853}{250000}$ & $\tfrac{1338267981}{40000}$ & $\tfrac{3072760897}{1000000000000000}$ & 31426/17 & 59 & 6.4 \\
29 & $\tfrac{1338267981}{40000}$ & $\tfrac{9481137173}{250000}$ & $\tfrac{2672852847}{1000000000000000}$ & 35690/18 & 64 & 7.1 \\
30 & $\tfrac{9481137173}{250000}$ & $\tfrac{1070418171}{25000}$ & $\tfrac{2335546837}{1000000000000000}$ & 40370/18 & 70 & 8.2 \\
31 & $\tfrac{1070418171}{25000}$ & $\tfrac{48110318623}{1000000}$ & $\tfrac{2050940953}{1000000000000000}$ & 45463/18 & 76 & 9.6 \\
32 & $\tfrac{48110318623}{1000000}$ & $\tfrac{10789326747}{200000}$ & $\tfrac{1807809003}{1000000000000000}$ & 51028/18 & 82 & 10 \\
33 & $\tfrac{10789326747}{200000}$ & $\tfrac{30165067587}{500000}$ & $\tfrac{1597827091}{1000000000000000}$ & 57137/18 & 89 & 12 \\
34 & $\tfrac{30165067587}{500000}$ & $\tfrac{33566442549}{500000}$ & $\tfrac{1418388113}{1000000000000000}$ & 63731/18 & 97 & 15 \\
35 & $\tfrac{33566442549}{500000}$ & $\tfrac{37283397387}{500000}$ & $\tfrac{157973183}{125000000000000}$ & 70849/19 & 105 & 19 \\
\bottomrule
\end{longtable}
\endgroup

The endpoint list satisfies $U_i=L_{i+1}$ for every adjacent pair.  Two
declarations in the linked
\codelink{Robin1984/Finite/Certificates/AllRows.lean}{AllRows module}
complete the aggregate check:
\begin{itemize}[leftmargin=1.5em,itemsep=.15em]
  \item \path{robinFiniteRows_cover} checks the gap-free equalities and outer
    endpoints by kernel reduction;
  \item \path{robinFiniteRows_checks} dispatches every row to its separately
    checked predicate.
\end{itemize}
The exact remaining cutoffs, gain bound, height mantissa and exponent, and
endpoint logarithm bounds are stored in the linked
\dirlink{Robin1984/Finite/Certificates}{certificate directory}.

\section{Notation map}\label{app:notation}

\begingroup
\small
\renewcommand{\arraystretch}{1.12}
\begin{longtable}{@{}>{\raggedright\arraybackslash}p{.19\textwidth}
  >{\raggedright\arraybackslash}p{.57\textwidth}
  >{\raggedright\arraybackslash}p{.16\textwidth}@{}}
\toprule
symbol & meaning & first definition \\
\midrule
\endfirsthead
\toprule
symbol & meaning & first definition \\
\midrule
\endhead
$\sig(n)$
& Sum of the positive divisors of $n$.
& Section~\ref{sec:informal} \\
$\A(n)$
& Abundancy ratio $\sig(n)/n$.
& Section~\ref{sec:informal} \\
$Q_\eps(n)$
& Penalized abundancy $\A(n)n^{-\eps}$; its global maximizers are
  colossally abundant at parameter $\eps$.
& \eqref{eq:ca-objective} \\
$\M(n)$
& Logarithmic Robin margin; Robin's inequality is $\M(n)<0$.
& \eqref{eq:margin} \\
$g_{p,j},\tau_{p,j}$
& Marginal log gain for raising $p^{j-1}$ to $p^j$, and its gain per unit
  log height.
& \eqref{eq:event} \\
$\lambda(n)$
& Tangent slope $1/(\log n,\log\log n)$.
& \eqref{eq:lambda} \\
$\eps_0$
& Exact parameter $\log(12/11)/\log11$ tying $5040$ and $55440$.
& \eqref{eq:epsilon-zero} \\
$\mathcal B(c)$
& Finite prime-layer box determined by the cutoff vector $c$.
& Section~\ref{sec:finite} \\
$L_N,U_N$
& Rational lower and upper bounds for the logarithm.
& \eqref{eq:log-series} \\
$H$
& Logarithmic height $\log n$ in the analytic argument.
& Section~\ref{sec:rh-forward} \\
$S(H)$
& Natural RH comparison scale $H^{-1/2}/\log H$.
& \eqref{eq:natural-scale} \\
$I_m,K_m,T_m,B_m$
& Weighted prime error, one-zero kernel, explicit correction, and its
  RH majorant.
& \eqref{eq:weight}--\eqref{eq:Bmx} \\
$c_0$
& Multiplicity-counted inverse-square zero sum
  $\euler+2-\log(4\pi)$.
& \eqref{eq:Bmx} \\
$\E(H)$
& Uniform analytic debit in the complete exponent budget.
& \eqref{eq:mertens-scalar} \\
$\Pcost(H)$
& Nonnegative arithmetic cost forced by omitted primes or finite exponents.
& \eqref{eq:prime-cost} \\
$\vartheta,\psi$
& Chebyshev prime and prime-power counting functions.
& Section~\ref{sec:rh-forward} \\
$\Pprod(x),F(x)$
& Nicolas's finite Euler product and its normalized oscillation function.
& \eqref{eq:nicolas-function} \\
$k,K,J$
& Positive kernel and the weighted $\vartheta$- and $\psi$-error tails used
  in the Landau argument.
& \eqref{eq:KJ} \\
$L_P$
& Least common multiple of $1,\ldots,P$.
& Section~\ref{sec:lcm} \\
$D(P),T(P)$
& Log-height loss and finite prime-power reserve in the LCM transfer.
& \eqref{eq:height-loss}--\eqref{eq:tower-reserve} \\
\bottomrule
\end{longtable}
\endgroup

\par\bigskip
\begin{center}
\scshape Acknowledgements and attribution
\end{center}
\smallskip

The central mathematical result is Guy Robin's.  The oscillation input and
finite product $F$ are due to Jean-Louis Nicolas, and the positivity
principle used in that argument is due to Edmund Landau.  The colossally
abundant framework is due to Leonidas Alaoglu and Paul Erd\H{o}s.  Explicit
prime-function estimates used in the reconstruction are credited to J.
Barkley Rosser, Lowell Schoenfeld, and N. Costa Pereira, including Pereira's
corrigendum.  The formal reconstruction, certificate design, and exposition
are the work of the Robin1984 project; credits for external Lean
infrastructure are recorded above and in the repository metadata.

\begin{thebibliography}{99}

\bibitem{AlaogluErdos1944}
L. Alaoglu and P. Erd\H{o}s,
\emph{On highly composite and similar numbers},
Trans. Amer. Math. Soc. 56 (1944), 448--469.

\bibitem{CostaPereira1985}
N. Costa Pereira,
\emph{Estimates for the Chebyshev function $\psi(x)-\vartheta(x)$},
Math. Comp. 44 (1985), no. 169, 211--221,
\href{https://doi.org/10.2307/2007805}{doi:10.2307/2007805}.

\bibitem{CostaPereira1987}
N. Costa Pereira,
\emph{Corrigendum to ``Estimates for the Chebyshev function
$\psi(x)-\vartheta(x)$''},
Math. Comp. 48 (1987), no. 177, 447.

\bibitem{GomesKontorovich2020}
B. Gomes and A. Kontorovich,
\emph{Riemann hypothesis in Lean (with or without Mathlib)},
formal development, 2020,
\url{https://github.com/AlexKontorovich/Lean-RH}.

\bibitem{Landau1905}
E. Landau,
\emph{\"Uber einen Satz von Tschebyschef},
Math. Ann. 61 (1905), 527--550,
\url{https://eudml.org/doc/158244}.

\bibitem{Lean4}
L. de Moura and S. Ullrich,
\emph{The Lean 4 theorem prover and programming language},
in: Automated Deduction--CADE 28, Lecture Notes in Computer Science 12699,
Springer, 2021, 625--635.

\bibitem{LoefflerStoll2025}
D. Loeffler and M. Stoll,
\emph{Formalizing zeta and $L$-functions in Lean},
Ann. Formalized Math. 1 (2025),
\href{https://doi.org/10.46298/afm.15328}{doi:10.46298/afm.15328}.

\bibitem{Mathlib}
The Mathlib Community,
\emph{The Lean mathematical library},
CPP 2020, 367--381.

\bibitem{Nicolas1983}
J.-L. Nicolas,
\emph{Petites valeurs de la fonction d'Euler},
J. Number Theory 17 (1983), no. 3, 375--388,
\href{https://doi.org/10.1016/0022-314X(83)90055-0}
{doi:10.1016/0022-314X(83)90055-0}.

\bibitem{Nicolas2012}
J.-L. Nicolas,
\emph{Small values of the Euler function and the Riemann hypothesis},
Acta Arith. 155 (2012), no. 3, 311--321,
\href{https://doi.org/10.4064/aa155-3-7}{doi:10.4064/aa155-3-7}.

\bibitem{PrimeNumberTheoremAnd}
The PrimeNumberTheoremAnd contributors,
\emph{PrimeNumberTheoremAnd: a Lean formalization of the prime number theorem
and related results},
source repository,
\url{https://github.com/AlexKontorovich/PrimeNumberTheoremAnd}.

\bibitem{Robin1984}
G. Robin,
\emph{Grandes valeurs de la fonction somme des diviseurs et hypoth\`ese de
Riemann},
J. Math. Pures Appl. (9) 63 (1984), no. 2, 187--213, MR0774171
(in French).

\bibitem{RosserSchoenfeld1962}
J. B. Rosser and L. Schoenfeld,
\emph{Approximate formulas for some functions of prime numbers},
Illinois J. Math. 6 (1962), 64--94,
\href{https://doi.org/10.1215/ijm/1255631807}
{doi:10.1215/ijm/1255631807}.

\end{thebibliography}

\end{document}
